\documentclass[11pt,a4paper]{article}
% Preámbulo común de las resoluciones de TPs de IA.
% Uso, desde temaN/resolucion-tpM.tex:  % Preámbulo común de las resoluciones de TPs de IA.
% Uso, desde temaN/resolucion-tpM.tex:  % Preámbulo común de las resoluciones de TPs de IA.
% Uso, desde temaN/resolucion-tpM.tex:  \input{../plantilla/preambulo}
\usepackage[T1]{fontenc}
\usepackage[utf8]{inputenc}
\usepackage[spanish,es-noshorthands,es-tabla]{babel}
\usepackage{lmodern}
\usepackage{microtype}
\usepackage[margin=2.3cm]{geometry}
\usepackage{xcolor}
\usepackage{booktabs}
\usepackage{tabularx}
\usepackage{array}
\usepackage{enumitem}
\usepackage{tcolorbox}
\usepackage{titlesec}
\usepackage{amsmath,amssymb}
\usepackage{graphicx}
\usepackage{caption}
\usepackage{tikz}
\usepackage[hang,flushmargin]{footmisc}
\usepackage[hidelinks]{hyperref}

\usetikzlibrary{babel,arrows.meta,positioning,calc,fit,decorations.pathreplacing}
\tcbuselibrary{skins,breakable}

\definecolor{azul}{HTML}{1F4E79}
\definecolor{gris}{HTML}{6B7280}
\definecolor{ambar}{HTML}{B45309}
\definecolor{fondoazul}{HTML}{F3F7FB}
\definecolor{fondogris}{HTML}{F4F4F5}
\definecolor{fondoambar}{HTML}{FFF7ED}

\titleformat{\section}{\Large\bfseries\color{azul}}{\thesection}{0.6em}{}
\titleformat{\subsection}{\large\bfseries}{\thesubsection}{0.6em}{}

% Enunciado de la consigna (gris).
\newtcolorbox{enunciado}[1]{
  breakable, enhanced, colback=fondogris, colframe=gris,
  boxrule=0.4pt, left=6pt, right=6pt, top=4pt, bottom=4pt,
  title={\bfseries #1}, coltitle=white, colbacktitle=gris, fonttitle=\small, fontupper=\small}
% Ampliación de la respuesta (azul). No es una tcolorbox: las notas al pie dentro de
% una tcolorbox que se parte entre páginas se superponen con el contenido.
\newenvironment{extras}{%
  \par\medskip\noindent
  {\setlength{\fboxsep}{3pt}\colorbox{azul}{\color{white}\bfseries\small\ EXTRAS\ }}%
  \par\nobreak\vspace{1pt}{\color{azul}\hrule height 0.6pt}\nobreak\vspace{3pt}}
 {\par\vspace{3pt}{\color{azul}\hrule height 0.6pt}\medskip}
% Advertencias, ambigüedades del enunciado, erratas (ámbar).
\newtcolorbox{aviso}[1]{
  breakable, enhanced, colback=fondoambar, colframe=ambar,
  boxrule=0.5pt, left=6pt, right=6pt, top=5pt, bottom=5pt,
  title={\bfseries #1}, coltitle=white, colbacktitle=ambar, fonttitle=\small}

\setlist[itemize]{itemsep=2pt, topsep=3pt, parsep=0pt}
\setlist[enumerate]{itemsep=2pt, topsep=3pt, parsep=0pt}
\captionsetup{font=small, labelfont=bf}
\setlength{\parindent}{0pt}
\setlength{\parskip}{4pt}

% Citas textuales (en el idioma original, sin traducir).
\newcommand{\cita}[1]{``\textit{#1}''}

% Fuentes, siempre en nota al pie.
%   \catedra{Episodio I}{Agentes en Inteligencia Artificial}
%   \aima{25.1}{973}     \pmg{12.1}{443}
\newcommand{\catedra}[2]{\footnote{Cátedra, #1, «#2».}}
\newcommand{\aima}[2]{\footnote{AIMA, §#1, p.~#2.}}
\newcommand{\pmg}[2]{\footnote{PMG, §#1, p.~#2.}}
% Para varias fuentes en una sola nota: \fuentes{AIMA, §2.1, p.~34; PMG, §12.1, p.~443}
\newcommand{\fuentes}[1]{\footnote{#1.}}

% Aclaración inicial de abreviaturas.
\newcommand{\notafuentes}{%
  {\small\color{gris} Las fuentes van en notas al pie. \textbf{Cátedra}: diapositivas de A.~J.~García (2026).
  \textbf{AIMA}: Russell \& Norvig, \textit{Artificial Intelligence: A Modern Approach}, 3.ª ed. (2010).
  \textbf{PMG}: Poole, Mackworth \& Goebel, \textit{Computational Intelligence: A Logical Approach} (1998).\par}}

% Íconos reutilizables (pics de TikZ). Uso: \pic[scale=0.8] at (x,y) {robot};
\tikzset{
  pics/robot/.style={code={
    \draw[thick, fill=azul!15] (-0.35,0.15) rectangle (0.35,0.65);            % cabeza
    \fill[azul] (-0.15,0.42) circle (0.06) (0.15,0.42) circle (0.06);         % ojos
    \draw[thick] (0,0.65) -- (0,0.8); \fill[ambar] (0,0.83) circle (0.05);    % antena
    \draw[thick, fill=azul!15] (-0.45,-0.45) rectangle (0.45,0.12);           % cuerpo
    \draw[thick] (-0.45,-0.1) -- (-0.65,0.1) (0.45,-0.1) -- (0.65,0.1);       % brazos
    \draw[thick, fill=gris!40] (-0.25,-0.55) circle (0.12) (0.25,-0.55) circle (0.12); % ruedas
  }},
  pics/compu/.style={code={
    \draw[thick, fill=azul!15, rounded corners=1pt] (-0.55,-0.25) rectangle (0.55,0.5); % pantalla
    \node[font=\tiny\ttfamily\bfseries, text=azul] at (0,0.12) {</>};
    \draw[thick] (0,-0.25) -- (0,-0.4); \draw[thick] (-0.25,-0.4) -- (0.25,-0.4);       % pie
  }},
  pics/persona/.style={code={
    \draw[thick] (0,0.35) circle (0.1);
    \draw[thick] (0,0.25) -- (0,-0.05) (-0.15,0.15) -- (0.15,0.15)
                 (0,-0.05) -- (-0.12,-0.3) (0,-0.05) -- (0.12,-0.3);
  }},
  pics/arbol/.style={code={
    \draw[thick, brown!60!black] (0,-0.3) -- (0,0);
    \fill[green!50!black!70] (0,0.15) circle (0.2);
  }},
  pics/documento/.style={code={
    \draw[thick, fill=white] (-0.2,-0.28) -- (-0.2,0.28) -- (0.1,0.28) -- (0.2,0.18) -- (0.2,-0.28) -- cycle;
    \foreach \y in {0.1,0,-0.1} \draw[gris] (-0.12,\y) -- (0.12,\y);
  }},
  pics/sobre/.style={code={
    \draw[thick, fill=white] (-0.25,-0.17) rectangle (0.25,0.17);
    \draw[thick] (-0.25,0.17) -- (0,-0.02) -- (0.25,0.17);
  }},
}

\usepackage[T1]{fontenc}
\usepackage[utf8]{inputenc}
\usepackage[spanish,es-noshorthands,es-tabla]{babel}
\usepackage{lmodern}
\usepackage{microtype}
\usepackage[margin=2.3cm]{geometry}
\usepackage{xcolor}
\usepackage{booktabs}
\usepackage{tabularx}
\usepackage{array}
\usepackage{enumitem}
\usepackage{tcolorbox}
\usepackage{titlesec}
\usepackage{amsmath,amssymb}
\usepackage{graphicx}
\usepackage{caption}
\usepackage{tikz}
\usepackage[hang,flushmargin]{footmisc}
\usepackage[hidelinks]{hyperref}

\usetikzlibrary{babel,arrows.meta,positioning,calc,fit,decorations.pathreplacing}
\tcbuselibrary{skins,breakable}

\definecolor{azul}{HTML}{1F4E79}
\definecolor{gris}{HTML}{6B7280}
\definecolor{ambar}{HTML}{B45309}
\definecolor{fondoazul}{HTML}{F3F7FB}
\definecolor{fondogris}{HTML}{F4F4F5}
\definecolor{fondoambar}{HTML}{FFF7ED}

\titleformat{\section}{\Large\bfseries\color{azul}}{\thesection}{0.6em}{}
\titleformat{\subsection}{\large\bfseries}{\thesubsection}{0.6em}{}

% Enunciado de la consigna (gris).
\newtcolorbox{enunciado}[1]{
  breakable, enhanced, colback=fondogris, colframe=gris,
  boxrule=0.4pt, left=6pt, right=6pt, top=4pt, bottom=4pt,
  title={\bfseries #1}, coltitle=white, colbacktitle=gris, fonttitle=\small, fontupper=\small}
% Ampliación de la respuesta (azul). No es una tcolorbox: las notas al pie dentro de
% una tcolorbox que se parte entre páginas se superponen con el contenido.
\newenvironment{extras}{%
  \par\medskip\noindent
  {\setlength{\fboxsep}{3pt}\colorbox{azul}{\color{white}\bfseries\small\ EXTRAS\ }}%
  \par\nobreak\vspace{1pt}{\color{azul}\hrule height 0.6pt}\nobreak\vspace{3pt}}
 {\par\vspace{3pt}{\color{azul}\hrule height 0.6pt}\medskip}
% Advertencias, ambigüedades del enunciado, erratas (ámbar).
\newtcolorbox{aviso}[1]{
  breakable, enhanced, colback=fondoambar, colframe=ambar,
  boxrule=0.5pt, left=6pt, right=6pt, top=5pt, bottom=5pt,
  title={\bfseries #1}, coltitle=white, colbacktitle=ambar, fonttitle=\small}

\setlist[itemize]{itemsep=2pt, topsep=3pt, parsep=0pt}
\setlist[enumerate]{itemsep=2pt, topsep=3pt, parsep=0pt}
\captionsetup{font=small, labelfont=bf}
\setlength{\parindent}{0pt}
\setlength{\parskip}{4pt}

% Citas textuales (en el idioma original, sin traducir).
\newcommand{\cita}[1]{``\textit{#1}''}

% Fuentes, siempre en nota al pie.
%   \catedra{Episodio I}{Agentes en Inteligencia Artificial}
%   \aima{25.1}{973}     \pmg{12.1}{443}
\newcommand{\catedra}[2]{\footnote{Cátedra, #1, «#2».}}
\newcommand{\aima}[2]{\footnote{AIMA, §#1, p.~#2.}}
\newcommand{\pmg}[2]{\footnote{PMG, §#1, p.~#2.}}
% Para varias fuentes en una sola nota: \fuentes{AIMA, §2.1, p.~34; PMG, §12.1, p.~443}
\newcommand{\fuentes}[1]{\footnote{#1.}}

% Aclaración inicial de abreviaturas.
\newcommand{\notafuentes}{%
  {\small\color{gris} Las fuentes van en notas al pie. \textbf{Cátedra}: diapositivas de A.~J.~García (2026).
  \textbf{AIMA}: Russell \& Norvig, \textit{Artificial Intelligence: A Modern Approach}, 3.ª ed. (2010).
  \textbf{PMG}: Poole, Mackworth \& Goebel, \textit{Computational Intelligence: A Logical Approach} (1998).\par}}

% Íconos reutilizables (pics de TikZ). Uso: \pic[scale=0.8] at (x,y) {robot};
\tikzset{
  pics/robot/.style={code={
    \draw[thick, fill=azul!15] (-0.35,0.15) rectangle (0.35,0.65);            % cabeza
    \fill[azul] (-0.15,0.42) circle (0.06) (0.15,0.42) circle (0.06);         % ojos
    \draw[thick] (0,0.65) -- (0,0.8); \fill[ambar] (0,0.83) circle (0.05);    % antena
    \draw[thick, fill=azul!15] (-0.45,-0.45) rectangle (0.45,0.12);           % cuerpo
    \draw[thick] (-0.45,-0.1) -- (-0.65,0.1) (0.45,-0.1) -- (0.65,0.1);       % brazos
    \draw[thick, fill=gris!40] (-0.25,-0.55) circle (0.12) (0.25,-0.55) circle (0.12); % ruedas
  }},
  pics/compu/.style={code={
    \draw[thick, fill=azul!15, rounded corners=1pt] (-0.55,-0.25) rectangle (0.55,0.5); % pantalla
    \node[font=\tiny\ttfamily\bfseries, text=azul] at (0,0.12) {</>};
    \draw[thick] (0,-0.25) -- (0,-0.4); \draw[thick] (-0.25,-0.4) -- (0.25,-0.4);       % pie
  }},
  pics/persona/.style={code={
    \draw[thick] (0,0.35) circle (0.1);
    \draw[thick] (0,0.25) -- (0,-0.05) (-0.15,0.15) -- (0.15,0.15)
                 (0,-0.05) -- (-0.12,-0.3) (0,-0.05) -- (0.12,-0.3);
  }},
  pics/arbol/.style={code={
    \draw[thick, brown!60!black] (0,-0.3) -- (0,0);
    \fill[green!50!black!70] (0,0.15) circle (0.2);
  }},
  pics/documento/.style={code={
    \draw[thick, fill=white] (-0.2,-0.28) -- (-0.2,0.28) -- (0.1,0.28) -- (0.2,0.18) -- (0.2,-0.28) -- cycle;
    \foreach \y in {0.1,0,-0.1} \draw[gris] (-0.12,\y) -- (0.12,\y);
  }},
  pics/sobre/.style={code={
    \draw[thick, fill=white] (-0.25,-0.17) rectangle (0.25,0.17);
    \draw[thick] (-0.25,0.17) -- (0,-0.02) -- (0.25,0.17);
  }},
}

\usepackage[T1]{fontenc}
\usepackage[utf8]{inputenc}
\usepackage[spanish,es-noshorthands,es-tabla]{babel}
\usepackage{lmodern}
\usepackage{microtype}
\usepackage[margin=2.3cm]{geometry}
\usepackage{xcolor}
\usepackage{booktabs}
\usepackage{tabularx}
\usepackage{array}
\usepackage{enumitem}
\usepackage{tcolorbox}
\usepackage{titlesec}
\usepackage{amsmath,amssymb}
\usepackage{graphicx}
\usepackage{caption}
\usepackage{tikz}
\usepackage[hang,flushmargin]{footmisc}
\usepackage[hidelinks]{hyperref}

\usetikzlibrary{babel,arrows.meta,positioning,calc,fit,decorations.pathreplacing}
\tcbuselibrary{skins,breakable}

\definecolor{azul}{HTML}{1F4E79}
\definecolor{gris}{HTML}{6B7280}
\definecolor{ambar}{HTML}{B45309}
\definecolor{fondoazul}{HTML}{F3F7FB}
\definecolor{fondogris}{HTML}{F4F4F5}
\definecolor{fondoambar}{HTML}{FFF7ED}

\titleformat{\section}{\Large\bfseries\color{azul}}{\thesection}{0.6em}{}
\titleformat{\subsection}{\large\bfseries}{\thesubsection}{0.6em}{}

% Enunciado de la consigna (gris).
\newtcolorbox{enunciado}[1]{
  breakable, enhanced, colback=fondogris, colframe=gris,
  boxrule=0.4pt, left=6pt, right=6pt, top=4pt, bottom=4pt,
  title={\bfseries #1}, coltitle=white, colbacktitle=gris, fonttitle=\small, fontupper=\small}
% Ampliación de la respuesta (azul). No es una tcolorbox: las notas al pie dentro de
% una tcolorbox que se parte entre páginas se superponen con el contenido.
\newenvironment{extras}{%
  \par\medskip\noindent
  {\setlength{\fboxsep}{3pt}\colorbox{azul}{\color{white}\bfseries\small\ EXTRAS\ }}%
  \par\nobreak\vspace{1pt}{\color{azul}\hrule height 0.6pt}\nobreak\vspace{3pt}}
 {\par\vspace{3pt}{\color{azul}\hrule height 0.6pt}\medskip}
% Advertencias, ambigüedades del enunciado, erratas (ámbar).
\newtcolorbox{aviso}[1]{
  breakable, enhanced, colback=fondoambar, colframe=ambar,
  boxrule=0.5pt, left=6pt, right=6pt, top=5pt, bottom=5pt,
  title={\bfseries #1}, coltitle=white, colbacktitle=ambar, fonttitle=\small}

\setlist[itemize]{itemsep=2pt, topsep=3pt, parsep=0pt}
\setlist[enumerate]{itemsep=2pt, topsep=3pt, parsep=0pt}
\captionsetup{font=small, labelfont=bf}
\setlength{\parindent}{0pt}
\setlength{\parskip}{4pt}

% Citas textuales (en el idioma original, sin traducir).
\newcommand{\cita}[1]{``\textit{#1}''}

% Fuentes, siempre en nota al pie.
%   \catedra{Episodio I}{Agentes en Inteligencia Artificial}
%   \aima{25.1}{973}     \pmg{12.1}{443}
\newcommand{\catedra}[2]{\footnote{Cátedra, #1, «#2».}}
\newcommand{\aima}[2]{\footnote{AIMA, §#1, p.~#2.}}
\newcommand{\pmg}[2]{\footnote{PMG, §#1, p.~#2.}}
% Para varias fuentes en una sola nota: \fuentes{AIMA, §2.1, p.~34; PMG, §12.1, p.~443}
\newcommand{\fuentes}[1]{\footnote{#1.}}

% Aclaración inicial de abreviaturas.
\newcommand{\notafuentes}{%
  {\small\color{gris} Las fuentes van en notas al pie. \textbf{Cátedra}: diapositivas de A.~J.~García (2026).
  \textbf{AIMA}: Russell \& Norvig, \textit{Artificial Intelligence: A Modern Approach}, 3.ª ed. (2010).
  \textbf{PMG}: Poole, Mackworth \& Goebel, \textit{Computational Intelligence: A Logical Approach} (1998).\par}}

% Íconos reutilizables (pics de TikZ). Uso: \pic[scale=0.8] at (x,y) {robot};
\tikzset{
  pics/robot/.style={code={
    \draw[thick, fill=azul!15] (-0.35,0.15) rectangle (0.35,0.65);            % cabeza
    \fill[azul] (-0.15,0.42) circle (0.06) (0.15,0.42) circle (0.06);         % ojos
    \draw[thick] (0,0.65) -- (0,0.8); \fill[ambar] (0,0.83) circle (0.05);    % antena
    \draw[thick, fill=azul!15] (-0.45,-0.45) rectangle (0.45,0.12);           % cuerpo
    \draw[thick] (-0.45,-0.1) -- (-0.65,0.1) (0.45,-0.1) -- (0.65,0.1);       % brazos
    \draw[thick, fill=gris!40] (-0.25,-0.55) circle (0.12) (0.25,-0.55) circle (0.12); % ruedas
  }},
  pics/compu/.style={code={
    \draw[thick, fill=azul!15, rounded corners=1pt] (-0.55,-0.25) rectangle (0.55,0.5); % pantalla
    \node[font=\tiny\ttfamily\bfseries, text=azul] at (0,0.12) {</>};
    \draw[thick] (0,-0.25) -- (0,-0.4); \draw[thick] (-0.25,-0.4) -- (0.25,-0.4);       % pie
  }},
  pics/persona/.style={code={
    \draw[thick] (0,0.35) circle (0.1);
    \draw[thick] (0,0.25) -- (0,-0.05) (-0.15,0.15) -- (0.15,0.15)
                 (0,-0.05) -- (-0.12,-0.3) (0,-0.05) -- (0.12,-0.3);
  }},
  pics/arbol/.style={code={
    \draw[thick, brown!60!black] (0,-0.3) -- (0,0);
    \fill[green!50!black!70] (0,0.15) circle (0.2);
  }},
  pics/documento/.style={code={
    \draw[thick, fill=white] (-0.2,-0.28) -- (-0.2,0.28) -- (0.1,0.28) -- (0.2,0.18) -- (0.2,-0.28) -- cycle;
    \foreach \y in {0.1,0,-0.1} \draw[gris] (-0.12,\y) -- (0.12,\y);
  }},
  pics/sobre/.style={code={
    \draw[thick, fill=white] (-0.25,-0.17) rectangle (0.25,0.17);
    \draw[thick] (-0.25,0.17) -- (0,-0.02) -- (0.25,0.17);
  }},
}

\usepackage{needspace}

\title{\vspace{-1.5cm}\textbf{Trabajo Práctico 8 --- Resolución}\\[3pt]
\large Programación Lógica Rebatible (DeLP)}
\author{\normalsize Inteligencia Artificial --- UNS --- 2026}
\date{}

% Estilos TikZ comunes a este TP.
\tikzset{
  caja/.style={draw=gris, thick, rounded corners=3pt, fill=fondogris, align=center, font=\small, inner sep=4pt},
  cajaazul/.style={caja, fill=fondoazul, draw=azul},
  cajaambar/.style={caja, fill=fondoambar, draw=ambar},
  flecha/.style={-{Stealth[length=2.2mm]}, thick, draw=azul},
  % Reglas en los grafos: estricta (llena) y rebatible (punteada).
  estricta/.style={-{Stealth[length=2.2mm]}, thick, draw=azul},
  rebatible/.style={-{Stealth[length=2.2mm]}, thick, dashed, draw=ambar},
  conflicto/.style={{Stealth[length=2mm]}-{Stealth[length=2mm]}, thick, draw=red!70!black},
  lit/.style={draw=gris, rounded corners=2pt, fill=white, font=\small\ttfamily, inner sep=3pt},
}

% Notación DeLP.
\newcommand{\nf}{\textasciitilde}                 % negación fuerte en texto
\newcommand{\reb}{\mathrel{-\!\!\!\prec}}         % regla rebatible  L0 -< L1,...,Ln
\newcommand{\est}{\leftarrow}                      % regla estricta   L0 <- L1,...,Ln
% Regla en el texto: \regla{cabeza}{flecha}{cuerpo}
\newcommand{\regla}[3]{\text{\texttt{#1}}~\ensuremath{#2}~\text{\texttt{#3}}}
% Programa P2 (ejercicios 2 y 3).
\newcommand{\programaPdos}{%
\[
\Pi_2 = \left\{\begin{array}{l}
\mathit{asustada(tina)} \\
\mathit{gallina(tina)} \\
\mathit{gallina(peluche)} \\
\sim\mathit{vuela(peluche)} \\
\mathit{ave(X)} \est \mathit{gallina(X)}
\end{array}\right\}
\qquad
\Delta_2 = \left\{\begin{array}{l}
\mathit{vuela(X)} \reb \mathit{ave(X)} \\
\sim\mathit{vuela(X)} \reb \mathit{gallina(X)} \\
\mathit{vuela(X)} \reb \mathit{gallina(X), asustada(X)} \\
\mathit{anida\_en\_ramas\_altas(X)} \reb \mathit{vuela(X)}
\end{array}\right\}
\]}
% Reglas de P2 con nombre (ejercicios 2 y 3).
\newcommand{\reglasPdos}{%
\begin{center}
\small
\begin{tabular}{@{}ll@{\qquad}ll@{}}
E1: & \regla{ave(X)}{\est}{gallina(X)} & R2: & \regla{\nf vuela(X)}{\reb}{gallina(X)} \\
R1: & \regla{vuela(X)}{\reb}{ave(X)} & R3: & \regla{vuela(X)}{\reb}{gallina(X), asustada(X)} \\
 & & R4: & \regla{anida\_en\_ramas\_altas(X)}{\reb}{vuela(X)} \\
\end{tabular}
\end{center}}
\newcommand{\fuentesweb}[1]{\footnote{\raggedright #1.}}

\begin{document}
\maketitle
\vspace{-1cm}
\notafuentes

\indice

%======================================================================
\section{Ejercicio 1: Conceptos básicos}

%----------------------------------------------------------------------
\subsection{a)}

\begin{enunciado}{Consigna}
¿Qué es un programa lógico rebatible?
\end{enunciado}

Un \textbf{programa lógico rebatible} (PLR) es un par $(\Pi, \Delta)$ (o el conjunto $\Pi \cup \Delta$) formado por:\fuentes{Cátedra, Episodio X, parte 2, «Reglas rebatibles (defeasible rules)» y «Hechos (facts)»}
\begin{itemize}
  \item $\Pi$: un conjunto de \textbf{hechos} y \textbf{reglas estrictas}.
  \begin{itemize}
    \item \textit{Hecho}: un literal fijo (sin variables).
    \item \textit{Regla estricta} $L_0 \est L_1, \dots, L_n$: información \textbf{segura}.
  \end{itemize}
  \item $\Delta$: un conjunto de \textbf{reglas rebatibles} $L_0 \reb L_1, \dots, L_n$: información \textbf{tentativa}.
\end{itemize}
Un \textbf{literal} es un átomo $p(T_1,\dots,T_n)$ o un átomo negado ${\sim}p(T_1,\dots,T_n)$, donde ``\nf'' es la \textbf{negación fuerte}.

\begin{figure}[h]
\centering
\begin{tikzpicture}
  \node[cajaazul, minimum width=4.6cm, minimum height=2.3cm, text width=4.3cm] (pi) at (0,0)
    {\textbf{$\Pi$}: información segura\\[3pt]
     \footnotesize hechos: \texttt{de\_noche.}\\
     \footnotesize estrictas: \regla{\nf encendida}{\est}{apagada}};
  \node[cajaambar, minimum width=4.6cm, minimum height=2.3cm, text width=4.3cm] (de) at (5.6,0)
    {\textbf{$\Delta$}: información tentativa\\[3pt]
     \footnotesize rebatibles: \regla{\nf hay\_luz}{\reb}{de\_noche}};
  \node[font=\Large] at (2.8,0) {$\cup$};
  \draw[decorate, decoration={brace, amplitude=6pt, mirror}, thick, gris] (-2.3,-1.3) -- (7.9,-1.3)
    node[midway, below=7pt, font=\small\bfseries, text=azul] {programa lógico rebatible $(\Pi, \Delta)$};
\end{tikzpicture}
\caption{Componentes de un PLR (con reglas del inciso d como ejemplo).}
\end{figure}

\begin{extras}
\begin{itemize}
  \item DeLP provee un \textbf{lenguaje de representación de conocimiento} que permite representar información \textbf{tentativa} y \textbf{potencialmente en conflicto}. Ante una consulta, su mecanismo de inferencia construye y evalúa \textbf{todos} los argumentos relacionados con ella (a favor y en contra).\catedra{Episodio X, parte 2}{Defeasible Logic Programing (DeLP)}
  \item \textbf{Variables}: las mayúsculas en los literales denotan variables; una regla con variables es un \textbf{esquema de reglas}. Los hechos, en cambio, son los encargados de \textbf{propagar las instancias} en las derivaciones de las conclusiones.\fuentes{Cátedra, Episodio X, parte 2, «Reglas rebatibles (defeasible rules)» y «Hechos (facts)»}
  \item La base es la tesis doctoral de A.~J.~García, \textit{Programación Lógica Rebatible} (UNS, 2000), cap.~2, y el artículo de García y Simari, \textit{Defeasible logic programming: An argumentative approach} (TPLP, 2004).\catedra{Episodio X, parte 2}{Defeasible Logic Programing (DeLP)}
\end{itemize}
\end{extras}

%----------------------------------------------------------------------
\subsection{b)}

\begin{enunciado}{Consigna}
¿Cuál es la utilidad de contar con dos tipos de reglas: estrictas y rebatibles? ¿Qué tipo de información permite modelar cada una de estas reglas? ¿Cómo se distinguen estas reglas en el lenguaje?
\end{enunciado}

\medskip\noindent\textbf{Qué modela cada una.}\fuentes{Cátedra, Episodio X, parte 2, «Reglas rebatibles (defeasible rules)» y «Reglas estrictas (strict rules)»}
\begin{itemize}
  \item \textbf{Estricta}: información \textbf{segura} (no rebatible). La conexión entre el antecedente y la conclusión \textbf{se cumple siempre}: no puede haber excepciones. Ej.: \regla{\nf abierto}{\est}{cerrado}.
  \item \textbf{Rebatible}: información \textbf{tentativa}. La conexión entre el antecedente (las razones) y la conclusión se cumple en muchos casos, pero \textbf{puede haber excepciones}. Ej.: \regla{\nf cruzar}{\reb}{semaforo(rojo)}.
\end{itemize}

\paragraph{Utilidad de tener las dos.}
\begin{itemize}
  \item Con las rebatibles se pueden expresar \textbf{reglas generales} que se cumplen en la mayoría de los casos y, además, reglas que dicen \textbf{en qué circunstancias hay excepciones}; también \textbf{posturas u opiniones diferentes} que pueden venir de distintas fuentes.\catedra{Episodio X, parte 2}{Defeasible Logic Programing (DeLP)}
  \item Con las estrictas se representa lo que no admite excepciones. Si hay excepciones, no es estricta y debe representarse como rebatible; la decisión depende del \textbf{dominio de aplicación} del SBC y es responsabilidad de quien codifica el conocimiento.\catedra{Episodio X, parte 2}{Reglas estrictas (strict rules)}
\end{itemize}

\paragraph{Cómo se distinguen.} Sintácticamente, \textbf{solo por la flecha} (aunque la diferencia semántica es enorme):\catedra{Episodio X, parte 2}{Reglas estrictas (strict rules)}
\begin{center}
\small
\begin{tabularx}{\linewidth}{@{}l c X l@{}}
\toprule
\textbf{Regla} & \textbf{Notación} & \textbf{Información que modela} & \textbf{Ejemplo} \\
\midrule
Estricta & $L_0 \est L_1, \dots, L_n$ & segura: se cumple siempre, sin excepciones & \regla{\nf abierto}{\est}{cerrado} \\
Rebatible & $L_0 \reb L_1, \dots, L_n$ & tentativa: se cumple en muchos casos, puede tener excepciones & \regla{\nf cruzar}{\reb}{semaforo(rojo)} \\
\bottomrule
\end{tabularx}
\end{center}

\begin{extras}
\begin{itemize}
  \item \textbf{Lectura de una regla rebatible} \texttt{conclusión} $\reb$ \texttt{antecedente}: ``si hay razones para creer en el antecedente, entonces hay una razón (rebatible) para creer en la conclusión''.\catedra{Episodio X, parte 2}{Reglas rebatibles (defeasible rules)}
  \item \textbf{Cadena de excepciones}: cada regla puede ser excepción de la anterior. Según los hechos que perciba el agente, pueden derivar conclusiones las tres, solo dos, solo una o ninguna:\catedra{Episodio X, parte 2}{Reglas rebatibles (defeasible rules)}
  \begin{center}
  \small
  \begin{tabular}{@{}ll@{}}
    \regla{\nf cruzar}{\reb}{semaforo(rojo).} & regla general \\
    \regla{cruzar}{\reb}{semaforo(rojo), emergencia.} & excepción a la primera \\
    \regla{\nf cruzar}{\reb}{semaforo(rojo), emergencia, hay\_peligro.} & excepción a la segunda \\
  \end{tabular}
  \end{center}
\end{itemize}
\end{extras}

%----------------------------------------------------------------------
\subsection{c)}

\begin{enunciado}{Consigna}
¿Qué ventajas tiene usar DeLP como mecanismo de inferencia y/o como herramienta de representación de conocimiento frente a otros enfoques, como el de la lógica clásica o frente a Prolog?
\end{enunciado}

\medskip\noindent\textbf{Ventajas de DeLP.}\fuentes{Cátedra, Episodio X, parte 2, «Defeasible Logic Programming». Mismo resumen en Episodio X (tercera parte), «Comentarios finales»}
\begin{itemize}
  \item \textit{Representación}: permite representar información \textbf{tentativa} y información \textbf{potencialmente en conflicto}.
  \item \textit{Inferencia}:
  \begin{itemize}
    \item las conclusiones están \textbf{sustentadas con argumentos};
    \item las respuestas se basan en un \textbf{análisis exhaustivo} de todos los argumentos relacionados (a favor y en contra);
    \item el criterio de preferencia entre argumentos se puede \textbf{adaptar al dominio} de aplicación;
    \item se puede dar una \textbf{explicación detallada} de las respuestas.
  \end{itemize}
\end{itemize}

\paragraph{Frente a Prolog.}
\begin{itemize}
  \item Prolog estándar \textbf{no tiene negación fuerte}, y tampoco un mecanismo para \textbf{manejar información contradictoria}: si hay derivaciones para \texttt{comprar\_acciones(acme)} y para \texttt{no\_comprar\_acciones(acme)}, responde \texttt{true} a las dos.\catedra{Episodio VIII}{Extender el mecanismo de inferencia}
  \item Solo tiene \textbf{negación por falla} (\texttt{\textbackslash+}), que no puede ir en la cabeza de una regla ni en un hecho.\catedra{Episodio VIII}{Negación por falla en Prolog}
  \item DeLP agrega la negación fuerte (\nf) y las reglas rebatibles, y ante conclusiones contradictorias \textbf{decide cuál prevalece} comparando los argumentos a favor y en contra.\catedra{Episodio X, parte 2}{Defeasible Logic Programing (DeLP)}
\end{itemize}

\paragraph{Frente a la lógica clásica.}
\begin{itemize}
  \item La lógica clásica es \textbf{monotónica}:\aima{7.5}{251} \cita{the set of entailed sentences can only increase as information is added to the knowledge base}. Lo que se concluyó no se retracta al aparecer una excepción.
  \item Para no inferir conclusiones contradictorias con inferencia monotónica (como Modus Ponens) hay que \textbf{mantener la base libre de contradicciones}. La otra opción es \textbf{admitir contradicciones potenciales} en la base y tener una inferencia (no monotónica) que \textbf{decida} ante una contradicción.\catedra{Episodio IX (valores de certeza)}{Bases de creencias inconsistentes: dos grandes opciones} DeLP sigue esta segunda opción.
\end{itemize}

\begin{center}
\small
\begin{tabularx}{\linewidth}{@{}>{\bfseries}l X X X@{}}
\toprule
 & \textbf{Lógica clásica} & \textbf{Prolog} & \textbf{DeLP} \\
\midrule
Negación & --- & solo por falla (\texttt{\textbackslash+}) & fuerte (\nf) \\
Información tentativa & no (monotónica) & no & sí (reglas rebatibles) \\
Contradicciones & base libre de contradicciones & concluye las dos & decide cuál prevalece \\
\bottomrule
\end{tabularx}
\end{center}

\begin{extras}
\begin{itemize}
  \item \textbf{No monotonía en AIMA}: razonar por defecto, sacando conclusiones \cita{in the absence of any reason to doubt it} y retractándolas si aparece nueva evidencia, es habitual en el sentido común. Ese razonamiento \cita{is said to exhibit nonmonotonicity, because the set of beliefs does not grow monotonically over time as new evidence arrives}. Ejemplos de lógicas no monotónicas: \textit{circumscription} y \textit{default logic}.\aima{12.6.1}{458}
  \item \textbf{Un paso intermedio sobre Prolog}: agregarle el operador \nf{} y una inferencia ``cauta'' que responda \texttt{si} (deriva $Q$ pero no su complemento), \texttt{no} (deriva el complemento pero no $Q$), \texttt{indeciso} (deriva ambos) o \texttt{desconozco} (ninguno). Ante una contradicción no decide: se abstiene.\catedra{Episodio VIII}{Extendiendo el mecanismo de inferencia}
\end{itemize}
\end{extras}

%----------------------------------------------------------------------
\subsection{d)}

\begin{enunciado}{Consigna}
Escriba reglas estrictas o rebatibles, según considere más adecuado, para representar el siguiente conocimiento:
\begin{enumerate}[label=\roman*.]
  \item Si la linterna está apagada, entonces no está encendida.
  \item Si la linterna está encendida, entonces no está apagada.
  \item Si es de noche, generalmente no hay luz.
  \item Si es de noche y tengo una linterna, entonces puedo iluminar.
  \item Si la linterna no tiene batería, entonces no ilumina.
  \item Si es de noche y hay luna llena, entonces hay luz.
\end{enumerate}
\end{enunciado}

\begin{aviso}{Supuesto}
``Puedo iluminar'' (iv) y ``no ilumina'' (v) se representan con el mismo literal \texttt{ilumina} (y su complemento), y ``hay luz'' (iii y vi) con \texttt{hay\_luz}; si no, las reglas no entrarían en conflicto.
\end{aviso}

Criterio: si la regla puede tener excepciones, es rebatible; si se cumple siempre, estricta.\catedra{Episodio X, parte 2}{Reglas estrictas (strict rules)}

\begin{center}
\small
\begin{tabularx}{\linewidth}{@{}c l c X@{}}
\toprule
 & \textbf{Regla} & \textbf{Tipo} & \textbf{Por qué} \\
\midrule
i & \regla{\nf encendida}{\est}{apagada} & estricta & por definición: apagada y encendida son opuestos, sin excepciones \\
ii & \regla{\nf apagada}{\est}{encendida} & estricta & ídem \\
iii & \regla{\nf hay\_luz}{\reb}{de\_noche} & rebatible & ``generalmente'': tiene excepciones (vi) \\
iv & \regla{ilumina}{\reb}{de\_noche, tengo\_linterna} & rebatible & tiene excepciones: v \\
v & \regla{\nf ilumina}{\est}{\nf tiene\_bateria} & estricta & sin batería, la linterna no ilumina nunca \\
vi & \regla{hay\_luz}{\reb}{de\_noche, luna\_llena} & rebatible & es una excepción a iii, y a su vez puede tener excepciones (p.~ej. noche nublada) \\
\bottomrule
\end{tabularx}
\end{center}

\begin{itemize}
  \item iv no dice ``generalmente'', pero v es una excepción: tener linterna de noche no basta si no tiene batería. Si iv fuera estricta, con \texttt{de\_noche}, \texttt{tengo\_linterna} y \texttt{\nf tiene\_bateria} en $\Pi$ se derivarían \texttt{ilumina} y \texttt{\nf ilumina} solo con información segura.
  \item iii y vi entran en conflicto (\texttt{\nf hay\_luz} vs.\ \texttt{hay\_luz}) en una noche de luna llena: las dos son razones tentativas y será la comparación de argumentos la que decida.
\end{itemize}

\begin{center}
\begin{tikzpicture}
  \tikzset{lbl/.style={font=\scriptsize, fill=white, inner sep=1pt, pos=0.5}}
  % antecedentes
  \node[lit] (luna) at (-3.6,0) {luna\_llena};
  \node[lit] (noche) at (-0.6,0) {de\_noche};
  \node[lit] (linterna) at (2.4,0) {tengo\_linterna};
  \node[lit] (bat) at (5.4,0) {\nf tiene\_bateria};
  % conclusiones
  \node[lit] (luz) at (-3.6,2) {hay\_luz};
  \node[lit] (noluz) at (-0.9,2) {\nf hay\_luz};
  \node[lit] (ilu) at (2.1,2) {ilumina};
  \node[lit] (noilu) at (4.8,2) {\nf ilumina};
  % reglas
  \draw[rebatible] (noche) -- node[lbl] {iii} (noluz);
  \draw[rebatible] (luna) -- node[lbl] {vi} (luz);
  \draw[rebatible] (noche) -- (luz);
  \draw[rebatible] (linterna) -- node[lbl] {iv} (ilu);
  \draw[rebatible] (noche) -- (ilu);
  \draw[estricta] (bat) -- node[lbl] {v} (noilu);
  % conflictos
  \draw[conflicto] (luz) -- (noluz);
  \draw[conflicto] (ilu) -- (noilu);
  \node[font=\scriptsize, text=red!70!black] at ($(luz)!0.5!(noluz)+(0,0.4)$) {conflicto};
  \node[font=\scriptsize, text=red!70!black] at ($(ilu)!0.5!(noilu)+(0,0.4)$) {conflicto};
  % i y ii
  \node[lit] (enc) at (8,2) {encendida};
  \node[lit] (noapag) at (10.6,2) {\nf apagada};
  \node[lit] (apag) at (8,0) {apagada};
  \node[lit] (noenc) at (10.6,0) {\nf encendida};
  \draw[estricta] (enc) -- node[lbl, above] {ii} (noapag);
  \draw[estricta] (apag) -- node[lbl, above] {i} (noenc);
  % leyenda
  \draw[estricta] (-3.6,-1) -- ++(0.8,0) node[right, font=\scriptsize, text=black] {estricta ($\est$)};
  \draw[rebatible] (0,-1) -- ++(0.8,0) node[right, font=\scriptsize, text=black] {rebatible ($\reb$)};
  \draw[conflicto] (3.4,-1) -- ++(0.8,0) node[right, font=\scriptsize, text=black] {literales complementarios};
\end{tikzpicture}
\captionof{figure}{Las reglas del inciso d: de los antecedentes a la conclusión, y los conflictos entre conclusiones.}
\end{center}

\begin{extras}
\begin{itemize}
  \item Con v estricta, si \texttt{\nf tiene\_bateria} es un hecho, el conjunto $\{$iv$\}$ junto con $\Pi$ es contradictorio (deriva \texttt{ilumina} y \texttt{\nf ilumina}), así que no es un argumento para \texttt{ilumina}: no cumple la condición de consistencia.\catedra{Episodio X, parte 2}{Argumento}
\end{itemize}
\end{extras}

%======================================================================
\section{Ejercicio 2}

\begin{enunciado}{Consigna}
Considere el siguiente programa lógico rebatible $\mathcal{P}_2 = (\Pi_2, \Delta_2)$:
\programaPdos
\end{enunciado}

En lo que sigue, las reglas se nombran así:
\reglasPdos

%----------------------------------------------------------------------
\subsection{a)}

\begin{enunciado}{Consigna}
¿A qué se denomina derivación rebatible? Obtenga todas las derivaciones para los siguientes literales: vuela(tina), \nf vuela(tina), asustada(tina), vuela(peluche) y \nf vuela(peluche).
\end{enunciado}

\paragraph{Derivación rebatible.} Una \textbf{derivación rebatible} para un literal $L$ a partir de un PLR $(\Pi, \Delta)$ es una secuencia finita de literales fijos $L_1, L_2, \dots, L_n = L$ tal que, para cada $L_i$:\catedra{Episodio X, parte 2}{Derivaciones rebatibles (defeasible derivation)}
\begin{itemize}
  \item $L_i$ es un \textbf{hecho} de $\Pi$, o
  \item existe en $(\Pi, \Delta)$ una \textbf{instancia de una regla} (estricta o rebatible) con cabeza $L_i$ y cuerpo $B_1, \dots, B_k$, donde todo literal del cuerpo ya está en $L_1, \dots, L_{i-1}$.
\end{itemize}
Si solo usa hechos y reglas estrictas de $\Pi$, es una \textbf{derivación estricta}.

\begin{aviso}{Qué se cuenta como ``todas''}
Cambiar el orden de literales independientes (p.~ej.\ \texttt{asustada(tina)} antes o después de \texttt{gallina(tina)}) o agregar literales que no se usan da otras secuencias que también cumplen la definición. Se listan las derivaciones distintas en cuanto a \textbf{qué hechos y reglas usan}.
\end{aviso}

\begin{center}
\small
\begin{tabularx}{\linewidth}{@{}l l >{\ttfamily}l X l@{}}
\toprule
\textbf{Literal} & & \textbf{\normalfont Secuencia} & \textbf{Justificación de cada $L_i$} & \textbf{Tipo} \\
\midrule
\texttt{vuela(tina)} & D1 & gallina(tina), & hecho & rebatible \\
 & & ave(tina), & E1 instanciada & \\
 & & vuela(tina) & R1 instanciada & \\
\cmidrule(l){2-5}
 & D2 & gallina(tina), & hecho & rebatible \\
 & & asustada(tina), & hecho & \\
 & & vuela(tina) & R3 instanciada & \\
\midrule
\texttt{\nf vuela(tina)} & D3 & gallina(tina), & hecho & rebatible \\
 & & \nf vuela(tina) & R2 instanciada & \\
\midrule
\texttt{asustada(tina)} & D4 & asustada(tina) & hecho & estricta \\
\midrule
\texttt{vuela(peluche)} & D5 & gallina(peluche), & hecho & rebatible \\
 & & ave(peluche), & E1 instanciada & \\
 & & vuela(peluche) & R1 instanciada & \\
\midrule
\texttt{\nf vuela(peluche)} & D6 & \nf vuela(peluche) & hecho & estricta \\
\cmidrule(l){2-5}
 & D7 & gallina(peluche), & hecho & rebatible \\
 & & \nf vuela(peluche) & R2 instanciada & \\
\bottomrule
\end{tabularx}
\end{center}

\begin{itemize}
  \item R3 no sirve para \texttt{vuela(peluche)}: no hay forma de derivar \texttt{asustada(peluche)}.
  \item R4 no aparece porque ninguno de los literales pedidos es \texttt{anida\_en\_ramas\_altas(...)}.
\end{itemize}

\begin{center}
\begin{tikzpicture}[
    hecho/.style={lit, font=\scriptsize\ttfamily, fill=fondogris},
    concl/.style={lit, font=\scriptsize\ttfamily},
    rl/.style={font=\tiny, fill=white, inner sep=0.5pt},
    dn/.style={font=\scriptsize\bfseries, text=gris}]
  % D1
  \node[dn] at (0,3.2) {D1};
  \node[concl] (a) at (0,2.6) {vuela(tina)};
  \node[concl] (b) at (0,1.5) {ave(tina)};
  \node[hecho] (c) at (0,0.4) {gallina(tina)};
  \draw[rebatible] (b) -- node[rl] {R1} (a);
  \draw[estricta] (c) -- node[rl] {E1} (b);
  % D2
  \node[dn] at (3.2,3.2) {D2};
  \node[concl] (a) at (3.2,2.6) {vuela(tina)};
  \node[hecho] (b) at (2.4,1.1) {gallina(tina)};
  \node[hecho] (c) at (4.1,0.4) {asustada(tina)};
  \draw[rebatible] (b) -- (a);
  \draw[rebatible] (c) -- node[rl] {R3} (a);
  % D3
  \node[dn] at (6.3,3.2) {D3};
  \node[concl] (a) at (6.3,2.6) {\nf vuela(tina)};
  \node[hecho] (b) at (6.3,1.1) {gallina(tina)};
  \draw[rebatible] (b) -- node[rl] {R2} (a);
  % D4
  \node[dn] at (8.6,3.2) {D4};
  \node[hecho] at (8.6,2.6) {asustada(tina)};
  % D5
  \node[dn] at (11,3.2) {D5};
  \node[concl] (a) at (11,2.6) {vuela(peluche)};
  \node[concl] (b) at (11,1.5) {ave(peluche)};
  \node[hecho] (c) at (11,0.4) {gallina(peluche)};
  \draw[rebatible] (b) -- node[rl] {R1} (a);
  \draw[estricta] (c) -- node[rl] {E1} (b);
  % D6
  \node[dn] at (13.6,3.2) {D6};
  \node[hecho] at (13.6,2.6) {\nf vuela(peluche)};
  % D7
  \node[dn] at (13.6,1.9) {D7};
  \node[concl] (a) at (13.6,1.3) {\nf vuela(peluche)};
  \node[hecho] (b) at (13.6,0.2) {gallina(peluche)};
  \draw[rebatible] (b) -- node[rl] {R2} (a);
  % leyenda
  \node[hecho] at (0.6,-0.7) {hecho};
  \draw[estricta] (2,-0.7) -- ++(0.7,0) node[right, font=\scriptsize] {estricta};
  \draw[rebatible] (4.3,-0.7) -- ++(0.7,0) node[right, font=\scriptsize] {rebatible};
\end{tikzpicture}
\captionof{figure}{Las derivaciones D1--D7 como árboles: cada literal se obtiene de los de abajo con la regla indicada.}
\end{center}

\begin{extras}
\begin{itemize}
  \item D1 y D5 muestran que una derivación puede usar \textbf{los dos tipos de reglas} (E1 estricta y R1 rebatible).\catedra{Episodio X, parte 2}{Derivaciones para literales complementarios}
  \item Los literales complementarios respecto de ``\nf'' (como \texttt{vuela(tina)} y \texttt{\nf vuela(tina)}) también se llaman \textbf{literales contradictorios} o \textbf{en conflicto}, porque representan dos conclusiones contradictorias.\catedra{Episodio X, parte 2}{Derivaciones para literales complementarios}
\end{itemize}
\end{extras}

%----------------------------------------------------------------------
\subsection{b)}

\begin{enunciado}{Consigna}
¿Cuándo se dice que un conjunto de reglas es contradictorio? Brinde un ejemplo.
\end{enunciado}

Un conjunto $S$ de hechos, reglas estrictas y rebatibles es \textbf{contradictorio} si a partir de $S$ se pueden obtener derivaciones rebatibles tanto para un literal $p$ como para su complemento ${\sim}p$.\catedra{Episodio X, parte 2}{Conjunto de reglas contradictorio}

\textbf{Ejemplo} (de la cátedra):\catedra{Episodio X, parte 2}{Conjunto de reglas contradictorio}
\[
\Pi_c = \left\{\begin{array}{l}
\texttt{en\_supermercado.} \\
\texttt{en\_casa.} \\
\texttt{\nf en\_casa} \est \texttt{en\_supermercado.}
\end{array}\right\}
\]
\begin{itemize}
  \item \texttt{en\_casa} se deriva porque es un hecho.
  \item \texttt{\nf en\_casa} se deriva con la secuencia \texttt{en\_supermercado}, \texttt{\nf en\_casa}.
\end{itemize}

\begin{center}
\begin{tikzpicture}
  \node[lit, fill=fondogris] (s) at (0,0) {en\_supermercado};
  \node[lit] (nc) at (0,1.4) {\nf en\_casa};
  \node[lit, fill=fondogris] (c) at (3,1.4) {en\_casa};
  \draw[estricta] (s) -- (nc);
  \draw[conflicto] (nc) -- node[above, font=\scriptsize, text=red!70!black] {conflicto} (c);
\end{tikzpicture}
\captionof{figure}{$\Pi_c$ deriva dos literales complementarios: es contradictorio.}
\end{center}

%----------------------------------------------------------------------
\subsection{c)}

\begin{enunciado}{Consigna}
¿Cuál es la utilidad de que un programa lógico rebatible sea contradictorio? ¿Tiene sentido que el conjunto $\Pi$ sea contradictorio?
\end{enunciado}

\paragraph{Utilidad de un PLR contradictorio.} Es lo que permite representar:\catedra{Episodio X, parte 2}{Defeasible Logic Programing (DeLP)}
\begin{itemize}
  \item \textbf{información tentativa}: reglas generales y reglas para sus \textbf{excepciones}. En $\mathcal{P}_2$, R1 (las aves vuelan) tiene la excepción R2 (las gallinas no vuelan), que a su vez tiene la excepción R3 (las gallinas asustadas vuelan);
  \item \textbf{información potencialmente en conflicto}: distintas posturas u opiniones, que pueden venir de distintas fuentes.
\end{itemize}
Las contradicciones no se evitan: las resuelve el mecanismo de inferencia, analizando los argumentos a favor y en contra.

\paragraph{¿$\Pi$ contradictorio?} No. En todo PLR $(\Pi, \Delta)$ \textbf{se asume que $\Pi$ no es contradictorio}: los hechos y las reglas estrictas (información segura) no pueden inferir conclusiones contradictorias.\catedra{Episodio X, parte 2}{Conjunto de reglas contradictorio}
\begin{itemize}
  \item Un conjunto contradictorio como $\Pi_c$ del inciso b \textbf{no es válido} en DeLP como $\Pi$.
  \item Además, un argumento $A$ exige que $\Pi \cup A$ no sea contradictorio. Con $\Pi$ contradictorio, $\Pi \cup A$ lo sería para todo $A$, y no habría ningún argumento.\catedra{Episodio X, parte 2}{Argumento}
\end{itemize}

%----------------------------------------------------------------------
\subsection{d)}

\begin{enunciado}{Consigna}
Justifique por qué $(\Pi_2 \cup \Delta_2)$ es contradictorio, y por qué $\Pi_2$ no es contradictorio.
\end{enunciado}

\begin{itemize}
  \item \textbf{$\Pi_2 \cup \Delta_2$ es contradictorio}: deriva dos pares de literales complementarios (inciso a):
  \begin{itemize}
    \item \texttt{vuela(tina)} (D1 o D2) y \texttt{\nf vuela(tina)} (D3);
    \item \texttt{vuela(peluche)} (D5) y \texttt{\nf vuela(peluche)} (D6 o D7).
  \end{itemize}
  \item \textbf{$\Pi_2$ no es contradictorio}: solo con hechos y E1 se derivan los cuatro hechos más \texttt{ave(tina)} y \texttt{ave(peluche)} (E1 aplicada a cada gallina). Entre estos seis literales no hay ningún par $p$, ${\sim}p$: el único literal negado, \texttt{\nf vuela(peluche)}, no tiene su complemento \texttt{vuela(peluche)}, que solo se obtiene con R1.
\end{itemize}

\begin{center}
\small
\begin{tabular}{@{}l l l@{}}
\toprule
\textbf{Literal} & \textbf{Desde $\Pi_2$} & \textbf{Desde $\Pi_2 \cup \Delta_2$} \\
\midrule
\texttt{asustada(tina)}, \texttt{gallina(tina)}, \texttt{gallina(peluche)} & sí (hechos) & sí \\
\texttt{ave(tina)}, \texttt{ave(peluche)} & sí (E1) & sí \\
\texttt{\nf vuela(peluche)} & sí (hecho) & sí \\
\texttt{vuela(peluche)} & no & \textbf{sí} (D5) \quad $\Rightarrow$ conflicto \\
\texttt{vuela(tina)} & no & \textbf{sí} (D1, D2) \\
\texttt{\nf vuela(tina)} & no & \textbf{sí} (D3) \quad $\Rightarrow$ conflicto con la anterior \\
\bottomrule
\end{tabular}
\end{center}

\begin{center}
\begin{tikzpicture}[hecho/.style={lit, font=\scriptsize\ttfamily, fill=fondogris},
                    concl/.style={lit, font=\scriptsize\ttfamily},
                    rl/.style={font=\tiny, fill=white, inner sep=0.5pt}]
  % tina
  \node[hecho] (gt) at (0,0) {gallina(tina)};
  \node[hecho] (at) at (3,0) {asustada(tina)};
  \node[concl] (avt) at (-1.6,1.4) {ave(tina)};
  \node[concl] (vt) at (0,2.8) {vuela(tina)};
  \node[concl] (nvt) at (3,2.8) {\nf vuela(tina)};
  \draw[estricta] (gt) -- node[rl] {E1} (avt);
  \draw[rebatible] (avt) -- node[rl] {R1} (vt);
  \draw[rebatible] (gt) -- node[rl, pos=0.35] {R3} (vt);
  \draw[rebatible] (at) -- (vt);
  \draw[rebatible] (gt) -- node[rl, pos=0.3] {R2} (nvt);
  \draw[conflicto] (vt) -- (nvt);
  % peluche
  \node[hecho] (gp) at (7.5,0) {gallina(peluche)};
  \node[concl] (avp) at (6.3,1.4) {ave(peluche)};
  \node[concl] (vp) at (7.2,2.8) {vuela(peluche)};
  \node[hecho] (nvp) at (10.4,2.8) {\nf vuela(peluche)};
  \draw[estricta] (gp) -- node[rl] {E1} (avp);
  \draw[rebatible] (avp) -- node[rl] {R1} (vp);
  \draw[rebatible] (gp) -- node[rl] {R2} (nvp);
  \draw[conflicto] (vp) -- (nvp);
  % marco de Pi2
  \node[font=\scriptsize, text=gris, anchor=west] at (-2.6,-0.75) {fondo gris: hechos de $\Pi_2$};
  \draw[estricta] (2.4,-0.75) -- ++(0.7,0) node[right, font=\scriptsize] {estricta};
  \draw[rebatible] (4.6,-0.75) -- ++(0.7,0) node[right, font=\scriptsize] {rebatible};
  \draw[conflicto] (6.8,-0.75) -- ++(0.7,0) node[right, font=\scriptsize] {literales complementarios};
\end{tikzpicture}
\captionof{figure}{Lo que se deriva de $\Pi_2 \cup \Delta_2$ para \texttt{vuela}. Los dos conflictos aparecen recién al usar reglas rebatibles (líneas punteadas).}
\end{center}

%======================================================================
\section{Ejercicio 3: Argumentos}

%----------------------------------------------------------------------
\subsection{a)}

\begin{enunciado}{Consigna}
¿Cómo se define formalmente la noción de \textit{argumento} en DeLP? Explique cada una de las condiciones exigidas en la definición de argumento.
\end{enunciado}

Un \textbf{argumento} para una conclusión $h$ a partir de un PLR $(\Pi, \Delta)$ es un par $\langle A, h\rangle$, donde $h$ es un literal fijo y $A$ es un conjunto de \textbf{reglas rebatibles instanciadas} de $\Delta$ que cumple:\catedra{Episodio X, parte 2}{Argumento}
\begin{enumerate}
  \item \textbf{Derivación}: existe una derivación rebatible para $h$ a partir de $A \cup \Pi$.
  \item \textbf{Consistencia}: $A \cup \Pi$ no es contradictorio.
  \item \textbf{Minimalidad}: no existe un subconjunto propio de $A$ que cumpla (1) y (2).
\end{enumerate}

\paragraph{Qué pide cada condición.} Con los ejemplos de la cátedra, a partir de:\catedra{Episodio X, parte 2}{Algunos detalles sobre las condiciones de los argumentos}
\[
\Pi_n = \left\{\begin{array}{l}
\texttt{\nf en\_casa} \est \texttt{en\_supermercado.} \\
\texttt{en\_supermercado.} \\
\texttt{de\_noche.} \\
\texttt{festejo.}
\end{array}\right\}
\qquad
\Delta_n = \left\{\begin{array}{l}
\texttt{en\_casa} \reb \texttt{de\_noche.} \\
\texttt{preparo\_cena} \reb \texttt{de\_noche.} \\
\texttt{\nf preparo\_cena} \reb \texttt{festejo.} \\
\texttt{acuesto\_tarde} \reb \texttt{festejo.} \\
\texttt{en\_oficina} \reb \texttt{trabajando.}
\end{array}\right\}
\]
\begin{itemize}
  \item \textbf{(1) Derivación}: las reglas de $A$, junto con los hechos y reglas estrictas, tienen que permitir \textbf{llegar a $h$}. $Q = \{\regla{en\_oficina}{\reb}{trabajando}\}$ no es argumento para \texttt{en\_oficina}: \texttt{trabajando} no se deriva, así que $Q \cup \Pi_n$ no deriva \texttt{en\_oficina}.
  \item \textbf{(2) Consistencia}: las reglas de $A$ \textbf{no pueden contradecir} lo que se sabe con seguridad ni contradecirse entre sí. $R = \{\regla{en\_casa}{\reb}{de\_noche}\}$ no es argumento para \texttt{en\_casa}: $R \cup \Pi_n$ deriva \texttt{en\_casa} y \texttt{\nf en\_casa}.
  \item \textbf{(3) Minimalidad}: que \textbf{no haya reglas de más}; todas son necesarias para la derivación. $S = \{\regla{\nf preparo\_cena}{\reb}{festejo},\ \regla{acuesto\_tarde}{\reb}{festejo}\}$ no es argumento para \texttt{\nf preparo\_cena}, porque $B = \{\regla{\nf preparo\_cena}{\reb}{festejo}\} \subset S$ ya cumple (1) y (2). $\langle B, \texttt{\nf preparo\_cena}\rangle$ sí es argumento.
\end{itemize}

\begin{center}
\small
\begin{tabularx}{\linewidth}{@{}l l >{\raggedright\arraybackslash}X >{\raggedright\arraybackslash}X@{}}
\toprule
\textbf{Condición} & \textbf{Exige} & \textbf{Contraejemplo en $(\Pi_n, \Delta_n)$} & \textbf{Por qué falla} \\
\midrule
(1) Derivación & $A \cup \Pi$ deriva $h$ & $Q$ para \texttt{en\_oficina} & no deriva \texttt{trabajando} \\
(2) Consistencia & $A \cup \Pi$ no contradictorio & $R$ para \texttt{en\_casa} & deriva \texttt{en\_casa} y \texttt{\nf en\_casa} \\
(3) Minimalidad & ningún $A' \subset A$ cumple (1) y (2) & $S$ para \texttt{\nf preparo\_cena} & $B \subset S$ ya alcanza \\
\bottomrule
\end{tabularx}
\end{center}

\begin{extras}
\begin{itemize}
  \item \textbf{Observaciones} sobre la definición:\catedra{Episodio X, parte 2}{Argumento}
  \begin{itemize}
    \item el conjunto incluye \textbf{solo reglas rebatibles} (instanciadas); los hechos y reglas estrictas que se usan quedan en $\Pi$;
    \item la conclusión es \textbf{única} y está instanciada;
    \item se dice que $\langle A, h\rangle$ \textbf{sustenta} $h$;
    \item para una misma conclusión puede haber \textbf{varios argumentos}, cada uno una razón diferente;
    \item el argumento para un \textbf{hecho} es el \textbf{conjunto vacío}.
  \end{itemize}
  \item \textbf{Notación gráfica}: un argumento se dibuja como un \textbf{triángulo}, con la conclusión en el vértice superior y el nombre del argumento adentro. Un sub-argumento es un triángulo dentro de otro.\catedra{Episodio X (tercera parte)}{Notación gráfica de argumentos}
\end{itemize}
\end{extras}

%----------------------------------------------------------------------
\subsection{b)}

\begin{enunciado}{Consigna}
Considerando el programa $\mathcal{P}_2$ del ejercicio anterior, muestre los argumentos posibles para cada uno de los siguientes literales: vuela(tina), \nf vuela(tina), asustada(tina), anida\_en\_ramas\_altas(tina).
\end{enunciado}

Programa $\mathcal{P}_2$, con los nombres de reglas del ejercicio 2:
\programaPdos
\reglasPdos

\begin{center}
\footnotesize
\begin{tabularx}{\linewidth}{@{}l >{\raggedright\arraybackslash}X c@{}}
\toprule
\textbf{Conclusión} & \textbf{Argumento} & \textbf{Derivación} \\
\midrule
\texttt{vuela(tina)} & $A_1 = \{\regla{vuela(tina)}{\reb}{ave(tina)}\}$ & D1 \\
 & $A_2 = \{\regla{vuela(tina)}{\reb}{gallina(tina), asustada(tina)}\}$ & D2 \\
\midrule
\texttt{\nf vuela(tina)} & $B = \{\regla{\nf vuela(tina)}{\reb}{gallina(tina)}\}$ & D3 \\
\midrule
\texttt{asustada(tina)} & $\emptyset$ (es un hecho) & D4 \\
\midrule
\texttt{anida\_en\_ramas\_altas(tina)} & $C_1 = \{\regla{anida\_en\_ramas\_altas(tina)}{\reb}{vuela(tina)},$\newline $\phantom{C_1 = \{}\regla{vuela(tina)}{\reb}{ave(tina)}\}$ & D1 + R4 \\
 & $C_2 = \{\regla{anida\_en\_ramas\_altas(tina)}{\reb}{vuela(tina)},$\newline $\phantom{C_2 = \{}\regla{vuela(tina)}{\reb}{gallina(tina), asustada(tina)}\}$ & D2 + R4 \\
\bottomrule
\end{tabularx}
\end{center}

Por qué cumplen las tres condiciones:
\begin{itemize}
  \item \textbf{(1)}: cada uno contiene exactamente las reglas rebatibles de la derivación indicada (D1--D4 del ejercicio 2; para $C_1$ y $C_2$, esa derivación seguida de \texttt{anida\_en\_ramas\_altas(tina)} por R4).
  \item \textbf{(2)}: con $\Pi_2$, ninguno deriva un literal y su complemento. Con $\Pi_2$ solo se deriva un literal negado, \texttt{\nf vuela(peluche)}, y estos argumentos no derivan \texttt{vuela(peluche)}. Además, ninguno junta R2 con R1 o R3 para tina.
  \item \textbf{(3)}: sacando cualquier regla ya no se deriva la conclusión.
\end{itemize}

Conjuntos que \textbf{no} son argumentos (ejemplos):
\begin{itemize}
  \item $A_1 \cup A_2$ para \texttt{vuela(tina)}: falla (3), porque $A_1 \subset A_1 \cup A_2$ ya alcanza.
  \item $C_1 \cup B$ para \texttt{anida\_en\_ramas\_altas(tina)}: falla (2), porque deriva \texttt{vuela(tina)} y \texttt{\nf vuela(tina)} (y también (3)).
\end{itemize}

\begin{center}
\begin{tikzpicture}
  % \tri{x}{y}{media base}{altura}{nombre}{conclusión}{color}
  \newcommand{\tri}[7]{%
    \draw[thick, draw=#7, fill=#7!25] (#1-#3,#2) -- (#1,#2+#4) -- (#1+#3,#2) -- cycle;
    \node[font=\scriptsize\bfseries] at (#1,#2+0.3*#4) {#5};
    \node[font=\scriptsize\ttfamily, above] at (#1,#2+#4) {#6};}
  \tri{0}{0}{0.55}{1.2}{$A_1$}{v(t)}{ambar}
  \tri{1.6}{0}{0.55}{1.2}{$A_2$}{v(t)}{ambar}
  \tri{3.4}{0}{0.55}{1.2}{$B$}{\nf v(t)}{azul}
  \node[font=\scriptsize\ttfamily] at (4.9,1.45) {as(t)};
  \node[font=\scriptsize] at (4.9,0.6) {$\emptyset$};
  % C1 con A1 adentro
  \tri{6.8}{0}{0.95}{2}{}{an(t)}{ambar}
  \node[font=\scriptsize\bfseries] at (7.45,0.25) {$C_1$};
  \tri{6.6}{0.12}{0.45}{1}{$A_1$}{}{ambar}
  \node[font=\scriptsize\ttfamily] at (6.6,1.25) {v(t)};
  % C2 con A2 adentro
  \tri{9.2}{0}{0.95}{2}{}{an(t)}{ambar}
  \node[font=\scriptsize\bfseries] at (9.85,0.25) {$C_2$};
  \tri{9}{0.12}{0.45}{1}{$A_2$}{}{ambar}
  \node[font=\scriptsize\ttfamily] at (9,1.25) {v(t)};
\end{tikzpicture}
\captionof{figure}{Los argumentos de b en notación de triángulos (\texttt{v(t)} = \texttt{vuela(tina)}, \texttt{as(t)} = \texttt{asustada(tina)}, \texttt{an(t)} = \texttt{anida\_en\_ramas\_altas(tina)}). $A_1$ y $A_2$ son sub-argumentos de $C_1$ y $C_2$.}
\end{center}

\begin{extras}
\begin{itemize}
  \item Un argumento $\langle S, q\rangle$ es \textbf{sub-argumento} de $\langle X, h\rangle$ si $S \subseteq X$. Por eso $\langle A_1, \texttt{vuela(tina)}\rangle$ es sub-argumento de $\langle C_1, \texttt{anida\_en\_ramas\_altas(tina)}\rangle$ (y $A_2$ de $C_2$).\catedra{Episodio X (tercera parte)}{Relaciones entre argumentos: sub-argumento}
\end{itemize}
\end{extras}

%----------------------------------------------------------------------
\subsection{c)}

\begin{enunciado}{Consigna}
¿Qué relación existe entre las nociones de argumento y derivación?
\end{enunciado}

\begin{itemize}
  \item Todo argumento \textbf{necesita una derivación}: es la condición (1), una derivación rebatible de $h$ a partir de $A \cup \Pi$.\catedra{Episodio X, parte 2}{Argumento}
  \item El conjunto $A$ son las \textbf{reglas rebatibles instanciadas que usa esa derivación}. En el ejemplo de la cátedra, $N$ tiene las dos reglas rebatibles (R2 y R5) de la derivación de \texttt{\nf comprar\_acciones(acme)}; la regla estricta E3 y el hecho que usa la derivación quedan en $\Pi$.\catedra{Episodio X, parte 2}{Argumento}
  \item La recíproca no vale: \textbf{no toda derivación da un argumento}. Además de derivar, $A$ tiene que ser consistente con $\Pi$ (2) y mínimo (3). Ejemplo: \texttt{vuela(peluche)} (inciso d).
  \item Si la derivación es \textbf{estricta}, no usa reglas rebatibles y el argumento es $\emptyset$ (p.~ej.\ D4 para \texttt{asustada(tina)}).\catedra{Episodio X, parte 2}{Argumento}
\end{itemize}

\begin{center}
\begin{tikzpicture}
  \node[cajaazul, text width=3.8cm] (d) at (0,0) {\textbf{Derivación} de $h$\\\footnotesize secuencia de literales con hechos, estrictas y rebatibles};
  \node[cajaambar, text width=3.8cm] (a) at (8.5,0) {\textbf{Argumento} $\langle A, h\rangle$\\\footnotesize $A$ = sus reglas rebatibles};
  \draw[flecha] (d) -- node[above, font=\scriptsize, align=center] {solo si $A \cup \Pi$ es\\consistente (2) y $A$ mínimo (3)} (a);
  \draw[flecha, draw=ambar] (a.south) to[bend left=25] node[below, font=\scriptsize] {siempre tiene una (1)} (d.south);
\end{tikzpicture}
\captionof{figure}{Todo argumento se apoya en una derivación; no toda derivación da un argumento.}
\end{center}

%----------------------------------------------------------------------
\subsection{d)}

\begin{enunciado}{Consigna}
¿Por qué a partir del programa $\mathcal{P}_2$ existe una derivación para vuela(peluche) pero no existe un argumento para vuela(peluche)?
\end{enunciado}

\begin{itemize}
  \item \textbf{Derivación}: D5 = \texttt{gallina(peluche)}, \texttt{ave(peluche)} (E1), \texttt{vuela(peluche)} (R1). Una derivación solo pide que cada literal se obtenga de un hecho o una regla; no controla contradicciones.
  \item \textbf{Argumento}: el único candidato es $A = \{\regla{vuela(peluche)}{\reb}{ave(peluche)}\}$ (R3 no sirve: \texttt{asustada(peluche)} no se deriva). Pero \texttt{\nf vuela(peluche)} es un \textbf{hecho} de $\Pi_2$, así que $A \cup \Pi_2$ deriva \texttt{vuela(peluche)} y \texttt{\nf vuela(peluche)}: es contradictorio y \textbf{no cumple la condición (2)}.
  \item Agregar reglas no lo arregla: cualquier conjunto que derive \texttt{vuela(peluche)}, unido con $\Pi_2$, sigue derivando también el hecho \texttt{\nf vuela(peluche)}.
\end{itemize}

\begin{center}
\begin{tikzpicture}[hecho/.style={lit, font=\scriptsize\ttfamily, fill=fondogris},
                    concl/.style={lit, font=\scriptsize\ttfamily},
                    rl/.style={font=\tiny, fill=white, inner sep=0.5pt}]
  \node[hecho] (gp) at (0,0) {gallina(peluche)};
  \node[concl] (avp) at (0,1.2) {ave(peluche)};
  \node[concl] (vp) at (0,2.4) {vuela(peluche)};
  \node[hecho] (nvp) at (3.4,2.4) {\nf vuela(peluche)};
  \draw[estricta] (gp) -- node[rl] {E1} (avp);
  \draw[rebatible] (avp) -- node[rl] {R1} (vp);
  \draw[conflicto] (vp) -- node[above=5pt, font=\scriptsize, text=red!70!black] {conflicto} (nvp);
  \node[font=\scriptsize, text=gris, anchor=west] at (2.2,1.5) {hecho de $\Pi_2$: siempre está};
  \draw[dashed, ambar, thick, rounded corners] (-1.35,-0.5) rectangle (1.35,2.8);
  \node[font=\scriptsize, text=ambar, anchor=north] at (0,-0.55) {D5: derivación \checkmark};
  \node[font=\scriptsize, text=red!70!black, anchor=west] at (1.6,0.4) {$A \cup \Pi_2$ contradictorio $\Rightarrow$ argumento \texttimes};
\end{tikzpicture}
\captionof{figure}{\texttt{vuela(peluche)} tiene derivación, pero choca con el hecho \texttt{\nf vuela(peluche)}.}
\end{center}

%======================================================================
\section{Ejercicio 4: Ataques y derrotas}

% Triángulo de argumento: \tri{x}{y}{media base}{altura}{nombre}{conclusión}{color}
\newcommand{\tri}[7]{%
  \draw[thick, draw=#7, fill=#7!25] (#1-#3,#2) -- (#1,#2+#4) -- (#1+#3,#2) -- cycle;
  \node[font=\scriptsize\bfseries] at (#1,#2+0.3*#4) {#5};
  \node[font=\scriptsize\ttfamily, above] at (#1,#2+#4) {#6};}
\tikzset{
  ataque/.style={{Stealth[length=2mm]}-{Stealth[length=2mm]}, thick, dashed, draw=gris},
  dpropia/.style={-{Stealth[length=2.4mm]}, very thick, draw=red!70!black},
  dbloqueo/.style={{Stealth[length=2.4mm]}-{Stealth[length=2.4mm]}, very thick, draw=red!70!black},
}

%----------------------------------------------------------------------
\subsection{a)}

\begin{enunciado}{Consigna}
Defina formalmente las nociones de contra-argumento y derrotador (distinguiendo entre derrotador propio y derrotador por bloqueo).
\end{enunciado}

Definiciones, dado un PLR $(\Pi, \Delta)$:
\begin{itemize}
  \item \textbf{Desacuerdo}: $\langle X, f\rangle$ y $\langle Y, g\rangle$ están en desacuerdo (o en conflicto) si $\Pi \cup \{f, g\}$ es contradictorio.\catedra{Episodio X (tercera parte)}{Relaciones entre argumentos: desacuerdo}
  \item \textbf{Contra-argumento (ataque)}: $\langle X, p\rangle$ ataca (contra-argumenta) a $\langle Y, h\rangle$ si existe un sub-argumento $\langle S, q\rangle$ de $\langle Y, h\rangle$ tal que $\langle X, p\rangle$ y $\langle S, q\rangle$ están en desacuerdo.\catedra{Episodio X (tercera parte)}{Ataque (contra-argumento)}
  \begin{itemize}
    \item $\langle S, q\rangle$ se llama \textbf{sub-argumento de desacuerdo}; puede ser el propio $\langle Y, h\rangle$.
    \item $q$ es el \textbf{punto de ataque}.
  \end{itemize}
  \item \textbf{Derrotador}: $\langle X, p\rangle$ derrota a $\langle Y, h\rangle$ si lo ataca con sub-argumento de desacuerdo $\langle S, q\rangle$ y se cumple una de estas condiciones:\catedra{Episodio X (tercera parte)}{Derrota propia y derrota por bloqueo}
  \begin{itemize}
    \item[(a)] $\langle X, p\rangle$ es \textbf{preferido} a $\langle S, q\rangle$: es un \textbf{derrotador propio};
    \item[(b)] $\langle X, p\rangle$ \textbf{no es comparable} con $\langle S, q\rangle$: es un \textbf{derrotador por bloqueo}.
  \end{itemize}
  La derrota combina ataque y preferencia, y la preferencia se evalúa entre el atacante y el \textbf{sub-argumento de desacuerdo} (no el argumento atacado completo).
\end{itemize}

\begin{center}
\small
\begin{tabular}{@{}l l l@{}}
\toprule
\textbf{$\langle X, p\rangle$ ataca a $\langle Y, h\rangle$ en $\langle S, q\rangle$, y\dots} & \textbf{¿Derrota?} & \textbf{Tipo} \\
\midrule
$\langle X, p\rangle$ es preferido a $\langle S, q\rangle$ & sí & derrotador propio \\
$\langle X, p\rangle$ y $\langle S, q\rangle$ no son comparables & sí & derrotador por bloqueo \\
$\langle S, q\rangle$ es preferido a $\langle X, p\rangle$ & no & (solo ataque) \\
\bottomrule
\end{tabular}
\end{center}

\begin{center}
\begin{minipage}{\linewidth}
\centering
\begin{tikzpicture}
  % Ataque en un sub-argumento
  \tri{0}{0}{0.5}{1.1}{$X$}{p}{azul}
  \tri{2.6}{0}{1.2}{2.4}{}{h}{ambar}
  \node[font=\scriptsize\bfseries] at (3.35,0.25) {$Y$};
  \tri{2.4}{0.35}{0.45}{0.9}{$S$}{}{ambar}
  \node[font=\scriptsize\ttfamily, above] at (2.4,1.22) {q};
  \draw[ataque] (0.25,1.15) -- (2.2,1.45);
  \node[font=\scriptsize, align=center] at (1.3,-0.5) {ataque en el sub-argumento\\de desacuerdo $\langle S, q\rangle$};
  % Derrota propia
  \tri{5.6}{0}{0.5}{1.1}{$X$}{p}{azul}
  \tri{7.4}{0}{0.5}{1.1}{$S$}{q}{ambar}
  \draw[dpropia] (5.95,0.7) -- (7.05,0.7);
  \node[font=\scriptsize, align=center] at (6.5,-0.5) {derrota propia\\($X$ preferido a $S$)};
  % Derrota por bloqueo
  \tri{9.6}{0}{0.5}{1.1}{$X$}{p}{azul}
  \tri{11.4}{0}{0.5}{1.1}{$S$}{q}{ambar}
  \draw[dbloqueo] (9.95,0.7) -- (11.05,0.7);
  \node[font=\scriptsize, align=center] at (10.5,-0.5) {derrota por bloqueo\\(no comparables)};
\end{tikzpicture}
\captionof{figure}{Ataque (flecha gris punteada), derrota propia (flecha simple) y derrota por bloqueo (flecha doble), como en la notación de la cátedra.}
\end{minipage}
\end{center}

\begin{extras}
\begin{itemize}
  \item Todo ataque a un sub-argumento es también un ataque al súper-argumento, y si $R$ derrota a un sub-argumento $X$ de $N$, también derrota a $N$.\catedra{Episodio X (tercera parte)}{Relaciones entre argumentos}
  \item Si $R$ derrota a $N$, $R$ \textbf{defiende indirectamente} todo lo que $N$ ataca; y un $D$ que ataque a $R$ defiende indirectamente a $N$ (y ataca indirectamente lo que $N$ ataca).\catedra{Episodio X (tercera parte)}{Relaciones entre argumentos}
\end{itemize}
\end{extras}

%----------------------------------------------------------------------
\subsection{b)}

\begin{enunciado}{Consigna}
A partir de los argumentos formados en el ejercicio 3.b, indique los ataques que aparecen entre los argumentos para vuela(tina) y los argumentos para \nf vuela(tina).
\end{enunciado}

Argumentos del ejercicio 3.b (nombres de reglas del ejercicio 2):
\begin{itemize}
  \item $\langle A_1, \texttt{vuela(tina)}\rangle$, con $A_1 = \{\regla{vuela(tina)}{\reb}{ave(tina)}\}$ (R1);
  \item $\langle A_2, \texttt{vuela(tina)}\rangle$, con $A_2 = \{\regla{vuela(tina)}{\reb}{gallina(tina), asustada(tina)}\}$ (R3);
  \item $\langle B, \texttt{\nf vuela(tina)}\rangle$, con $B = \{\regla{\nf vuela(tina)}{\reb}{gallina(tina)}\}$ (R2).
\end{itemize}

Ataques:
\begin{itemize}
  \item $A_1$ y $B$ \textbf{se atacan mutuamente}: $\Pi_2 \cup \{\texttt{vuela(tina)}, \texttt{\nf vuela(tina)}\}$ es contradictorio.
  \item $A_2$ y $B$ \textbf{se atacan mutuamente}, por la misma razón.
  \item Cada argumento tiene una sola regla, así que su único sub-argumento es él mismo: el sub-argumento de desacuerdo es el \textbf{argumento completo} y el punto de ataque es la \textbf{conclusión}.
  \item $A_1$ y $A_2$ \textbf{no se atacan}: tienen la misma conclusión, no están en desacuerdo.
\end{itemize}

\begin{center}
\small
\begin{tabular}{@{}c c c c@{}}
\toprule
\textbf{Atacante} & \textbf{Atacado} & \textbf{Sub-argumento de desacuerdo} & \textbf{Punto de ataque} \\
\midrule
$\langle A_1, \texttt{vuela(tina)}\rangle$ & $\langle B, \texttt{\nf vuela(tina)}\rangle$ & $\langle B, \texttt{\nf vuela(tina)}\rangle$ & \texttt{\nf vuela(tina)} \\
$\langle B, \texttt{\nf vuela(tina)}\rangle$ & $\langle A_1, \texttt{vuela(tina)}\rangle$ & $\langle A_1, \texttt{vuela(tina)}\rangle$ & \texttt{vuela(tina)} \\
$\langle A_2, \texttt{vuela(tina)}\rangle$ & $\langle B, \texttt{\nf vuela(tina)}\rangle$ & $\langle B, \texttt{\nf vuela(tina)}\rangle$ & \texttt{\nf vuela(tina)} \\
$\langle B, \texttt{\nf vuela(tina)}\rangle$ & $\langle A_2, \texttt{vuela(tina)}\rangle$ & $\langle A_2, \texttt{vuela(tina)}\rangle$ & \texttt{vuela(tina)} \\
\bottomrule
\end{tabular}
\end{center}

\begin{center}
\begin{minipage}{\linewidth}
\centering
\begin{tikzpicture}
  \tri{0}{0}{0.55}{1.2}{$A_1$}{v(t)}{ambar}
  \tri{3}{0}{0.55}{1.2}{$B$}{\nf v(t)}{azul}
  \tri{6}{0}{0.55}{1.2}{$A_2$}{v(t)}{ambar}
  \draw[ataque] (0.4,0.6) -- (2.6,0.6);
  \draw[ataque] (3.4,0.6) -- (5.6,0.6);
  \node[font=\scriptsize, text=gris] at (1.5,0.85) {ataque};
  \node[font=\scriptsize, text=gris] at (4.5,0.85) {ataque};
\end{tikzpicture}
\captionof{figure}{Ataques entre los argumentos para \texttt{vuela(tina)} (\texttt{v(t)}) y \texttt{\nf vuela(tina)}: $B$ ataca y es atacado por $A_1$ y por $A_2$.}
\end{minipage}
\end{center}

\begin{extras}
\begin{itemize}
  \item $B$ también ataca a $\langle C_1, \texttt{anida\_en\_ramas\_altas(tina)}\rangle$ y a $\langle C_2, \texttt{anida\_en\_ramas\_altas(tina)}\rangle$ del ejercicio 3.b, con sub-argumentos de desacuerdo $A_1$ y $A_2$ (punto de ataque \texttt{vuela(tina)}): todo ataque a un sub-argumento es un ataque al súper-argumento. En cambio, $C_1$ y $C_2$ no atacan a $B$: $\Pi_2 \cup \{\texttt{anida\_en\_ramas\_altas(tina)}, \texttt{\nf vuela(tina)}\}$ no es contradictorio.\catedra{Episodio X (tercera parte)}{Relaciones entre argumentos}
\end{itemize}
\end{extras}

%----------------------------------------------------------------------
\subsection{c)}

\begin{enunciado}{Consigna}
Defina un orden de preferencia para el conjunto de reglas rebatibles $\Delta_2$ y, a partir de las relaciones de ataque presentadas en el inciso anterior, indique qué derrotas se producen. Para cada una, indique si se trata de una derrota propia o por bloqueo.
\end{enunciado}

\paragraph{Orden elegido.} Se prioriza la regla más específica (la excepción sobre la regla general):
\[
\text{R2} > \text{R1} \quad (\text{\small las gallinas no vuelan, aunque sean aves}) \qquad\quad
\text{R3} > \text{R2} \quad (\text{\small las gallinas asustadas vuelan})
\]
Por transitividad, también R3 $>$ R1. R4 queda sin relacionar (no aparece en los argumentos que se comparan).

\paragraph{Criterio de preferencia (prioridad entre reglas).} $N$ es preferido a $A$ si:\catedra{Episodio X (tercera parte)}{Preferencia entre argumentos}
\begin{enumerate}
  \item existe $R_1 \in N$ y $R_2 \in A$ con $R_1 > R_2$ ($N$ tiene al menos una regla mejor que una de $A$), y
  \item no existe $R_1 \in N$ y $R_2 \in A$ con $R_1 < R_2$ ($A$ no tiene ninguna regla mejor que una de $N$).
\end{enumerate}

\paragraph{Comparaciones.}
\begin{itemize}
  \item $B$ vs.\ $A_1$: R2 $>$ R1 y $A_1$ no tiene reglas mejores $\Rightarrow$ \textbf{$B$ es preferido a $A_1$}.
  \item $A_2$ vs.\ $B$: R3 $>$ R2 y $B$ no tiene reglas mejores $\Rightarrow$ \textbf{$A_2$ es preferido a $B$}.
\end{itemize}

\paragraph{Derrotas.}
\begin{center}
\small
\begin{tabular}{@{}c c c c@{}}
\toprule
\textbf{Ataque} & \textbf{Preferencia} & \textbf{¿Derrota?} & \textbf{Tipo} \\
\midrule
$B$ ataca a $A_1$ & $B$ preferido a $A_1$ & \textbf{sí} & \textbf{propia} \\
$A_1$ ataca a $B$ & $B$ preferido a $A_1$ & no & --- \\
$A_2$ ataca a $B$ & $A_2$ preferido a $B$ & \textbf{sí} & \textbf{propia} \\
$B$ ataca a $A_2$ & $A_2$ preferido a $B$ & no & --- \\
\bottomrule
\end{tabular}
\end{center}
\begin{itemize}
  \item $\langle B, \texttt{\nf vuela(tina)}\rangle$ es \textbf{derrotador propio} de $\langle A_1, \texttt{vuela(tina)}\rangle$.
  \item $\langle A_2, \texttt{vuela(tina)}\rangle$ es \textbf{derrotador propio} de $\langle B, \texttt{\nf vuela(tina)}\rangle$.
  \item Con este orden \textbf{no hay derrotas por bloqueo}: todos los pares que se atacan son comparables.
\end{itemize}

\begin{center}
\begin{minipage}{\linewidth}
\centering
\begin{tikzpicture}
  \tri{0}{0}{0.55}{1.2}{$A_1$}{v(t)}{ambar}
  \tri{3}{0}{0.55}{1.2}{$B$}{\nf v(t)}{azul}
  \tri{6}{0}{0.55}{1.2}{$A_2$}{v(t)}{ambar}
  \draw[dpropia] (2.6,0.6) -- (0.4,0.6);
  \draw[dpropia] (5.6,0.6) -- (3.4,0.6);
  \node[font=\scriptsize, text=red!70!black] at (1.5,0.85) {propia};
  \node[font=\scriptsize, text=red!70!black] at (4.5,0.85) {propia};
  \node[font=\scriptsize] at (1.5,0.25) {R2 $>$ R1};
  \node[font=\scriptsize] at (4.5,0.25) {R3 $>$ R2};
\end{tikzpicture}
\captionof{figure}{Derrotas con el orden R3 $>$ R2 $>$ R1: $B$ derrota a $A_1$ y $A_2$ derrota a $B$, ambas derrotas propias.}
\end{minipage}
\end{center}

\begin{extras}
\begin{itemize}
  \item \textbf{Otro orden posible}: solo R2 $>$ R1 (sin comparar R3 con R2). Entonces $A_2$ y $B$ no son comparables y \textbf{se derrotan mutuamente por bloqueo}; $B$ sigue siendo derrotador propio de $A_1$. Es el mismo caso que $\langle A, g\rangle$ y $\langle C, {\sim}g\rangle$ en el ejemplo de la cátedra.\catedra{Episodio X (tercera parte)}{Derrota propia y derrota por bloqueo}
  \begin{center}
  \begin{tikzpicture}
    \tri{0}{0}{0.55}{1.2}{$A_1$}{v(t)}{ambar}
    \tri{3}{0}{0.55}{1.2}{$B$}{\nf v(t)}{azul}
    \tri{6}{0}{0.55}{1.2}{$A_2$}{v(t)}{ambar}
    \draw[dpropia] (2.6,0.6) -- (0.4,0.6);
    \draw[dbloqueo] (5.6,0.6) -- (3.4,0.6);
    \node[font=\scriptsize, text=red!70!black] at (1.5,0.85) {propia};
    \node[font=\scriptsize, text=red!70!black] at (4.5,0.85) {bloqueo};
  \end{tikzpicture}
  \end{center}
  \item Con el orden elegido, $A_2$ derrota a $B$, que derrota a $A_1$: $A_2$ \textbf{defiende indirectamente} a $A_1$.\catedra{Episodio X (tercera parte)}{Relaciones entre argumentos}
\end{itemize}
\end{extras}

%======================================================================
\section{Ejercicio 5}

% Programa P5 (ejercicios 5 y siguientes).
\newcommand{\programaPcinco}{%
\[
\Pi_5 = \left\{\begin{array}{l}
\mathit{viento} \est \mathit{viento\_caluroso} \\
\mathit{calor} \\
\mathit{caluroso} \\
\mathit{viento\_caluroso} \\
\mathit{ahorrar\_energia}
\end{array}\right\}
\qquad
\Delta_5 = \left\{\begin{array}{l}
\mathit{aire\_acond} \reb \mathit{calor} \\
\mathit{aire\_acond} \reb \mathit{sofocante} \\
\sim\mathit{aire\_acond} \reb \mathit{abrir\_ventana} \\
\sim\mathit{aire\_acond} \reb \mathit{ahorrar\_energia} \\
\mathit{abrir\_ventana} \reb \mathit{calor, viento} \\
\sim\mathit{abrir\_ventana} \reb \mathit{viento\_caluroso}
\end{array}\right\}
\]}
% Reglas de P5 con nombre.
\newcommand{\reglasPcinco}{%
\begin{center}
\small
\begin{tabular}{@{}ll@{\qquad}ll@{}}
R1: & \regla{aire\_acond}{\reb}{calor} & R4: & \regla{\nf aire\_acond}{\reb}{ahorrar\_energia} \\
R2: & \regla{aire\_acond}{\reb}{sofocante} & R5: & \regla{abrir\_ventana}{\reb}{calor, viento} \\
R3: & \regla{\nf aire\_acond}{\reb}{abrir\_ventana} & R6: & \regla{\nf abrir\_ventana}{\reb}{viento\_caluroso} \\
\end{tabular}
\end{center}}

\begin{enunciado}{Consigna}
Considere el programa $\mathcal{P}_5 = (\Pi_5, \Delta_5)$:
\programaPcinco
\end{enunciado}

En lo que sigue, las reglas se nombran así:
\reglasPcinco

\begin{aviso}{\texttt{sofocante} no se deriva}
R2 nunca se puede usar: \texttt{sofocante} no es un hecho ni cabeza de ninguna regla. En $\Pi_5$ está el hecho \texttt{caluroso}, que ninguna regla usa; probablemente R2 debía ser \regla{aire\_acond}{\reb}{caluroso}. Se resuelve el programa tal como está.
\end{aviso}

\clearpage
%----------------------------------------------------------------------
\subsection{a)}

\begin{enunciado}{Consigna}
Brinde dos argumentos para no encender el aire acondicionado (\nf aire\_acond) y muestre todos los contra-argumentos que los atacan.
\end{enunciado}

\paragraph{Argumentos para \texttt{\nf aire\_acond}.}
\begin{itemize}
  \item $\langle A_1, \texttt{\nf aire\_acond}\rangle$, con $A_1 = \{\regla{\nf aire\_acond}{\reb}{abrir\_ventana},\ \regla{abrir\_ventana}{\reb}{calor, viento}\}$ (R3, R5). \texttt{calor} es hecho y \texttt{viento} sale de la regla estricta con \texttt{viento\_caluroso}.
  \begin{itemize}
    \item Tiene el sub-argumento $\langle D, \texttt{abrir\_ventana}\rangle$, con $D = \{\regla{abrir\_ventana}{\reb}{calor, viento}\}$ (R5).
  \end{itemize}
  \item $\langle A_2, \texttt{\nf aire\_acond}\rangle$, con $A_2 = \{\regla{\nf aire\_acond}{\reb}{ahorrar\_energia}\}$ (R4).
\end{itemize}

\paragraph{Contra-argumentos.} Un contra-argumento tiene que estar en desacuerdo con algún sub-argumento: concluir \texttt{aire\_acond} (contra $A_1$ o $A_2$) o \texttt{\nf abrir\_ventana} (contra $D$).\catedra{Episodio X (tercera parte)}{Ataque (contra-argumento)}
\begin{itemize}
  \item $\langle B, \texttt{aire\_acond}\rangle$, con $B = \{\regla{aire\_acond}{\reb}{calor}\}$ (R1). Es el único argumento para \texttt{aire\_acond} (R2 no se puede usar).
  \item $\langle C, \texttt{\nf abrir\_ventana}\rangle$, con $C = \{\regla{\nf abrir\_ventana}{\reb}{viento\_caluroso}\}$ (R6).
\end{itemize}

\begin{center}
\small
\begin{tabular}{@{}c c c c@{}}
\toprule
\textbf{Contra-argumento} & \textbf{Ataca a} & \textbf{Sub-argumento de desacuerdo} & \textbf{Punto de ataque} \\
\midrule
$\langle B, \texttt{aire\_acond}\rangle$ & $\langle A_1, \texttt{\nf aire\_acond}\rangle$ & $\langle A_1, \texttt{\nf aire\_acond}\rangle$ & \texttt{\nf aire\_acond} \\
$\langle C, \texttt{\nf abrir\_ventana}\rangle$ & $\langle A_1, \texttt{\nf aire\_acond}\rangle$ & $\langle D, \texttt{abrir\_ventana}\rangle$ & \texttt{abrir\_ventana} \\
$\langle B, \texttt{aire\_acond}\rangle$ & $\langle A_2, \texttt{\nf aire\_acond}\rangle$ & $\langle A_2, \texttt{\nf aire\_acond}\rangle$ & \texttt{\nf aire\_acond} \\
\bottomrule
\end{tabular}
\end{center}

\begin{center}
\begin{minipage}{\linewidth}
\centering
\begin{tikzpicture}
  \tri{-3}{0}{0.55}{1.2}{$C$}{\nf ab}{azul}
  \tri{0}{0}{1.1}{2.4}{}{\nf ac}{ambar}
  \node[font=\scriptsize\bfseries] at (0.55,0.25) {$A_1$};
  \tri{-0.2}{0.35}{0.45}{0.9}{$D$}{}{ambar}
  \node[font=\scriptsize\ttfamily, above] at (-0.2,1.17) {ab};
  \tri{3.2}{0}{0.55}{1.2}{$B$}{ac}{azul}
  \tri{6.2}{0}{0.55}{1.2}{$A_2$}{\nf ac}{ambar}
  \draw[ataque] (-2.6,0.8) -- (-0.45,1.4);
  \draw[ataque] (2.85,0.9) -- (0.35,2.05);
  \draw[ataque] (3.6,0.6) -- (5.8,0.6);
\end{tikzpicture}
\captionof{figure}{Ataques a los argumentos para \texttt{\nf aire\_acond} (\texttt{ac} = \texttt{aire\_acond}, \texttt{ab} = \texttt{abrir\_ventana}). $C$ ataca a $A_1$ en su sub-argumento $D$.}
\end{minipage}
\end{center}

%----------------------------------------------------------------------
\subsection{b)}

\begin{enunciado}{Consigna}
Compare los argumentos obtenidos a partir del criterio de preferencia entre reglas visto en clase. En este caso, considere el siguiente orden:
\begin{center}
(\regla{aire\_acond}{\reb}{sofocante}) $>$ (\regla{\nf aire\_acond}{\reb}{ahorrar\_energia})\\
(\regla{\nf aire\_acond}{\reb}{ahorrar\_energia}) $>$ (\regla{aire\_acond}{\reb}{calor})\\
(\regla{\nf abrir\_ventana}{\reb}{viento\_caluroso}) $>$ (\regla{abrir\_ventana}{\reb}{calor, viento})
\end{center}
Determine si los contra-argumentos obtenidos en el inciso anterior son derrotadores. En caso afirmativo, indique si son derrotadores propios o por bloqueo.
\end{enunciado}

Con los nombres de reglas: \textbf{R2 $>$ R4}, \textbf{R4 $>$ R1}, \textbf{R6 $>$ R5} (y R2 $>$ R1 por transitividad). R3 no está relacionada con ninguna regla.

La preferencia se evalúa entre el atacante y el \textbf{sub-argumento de desacuerdo}; $N$ es preferido a $A$ si tiene una regla mejor que una de $A$ y $A$ no tiene ninguna mejor que una de $N$.\fuentes{Cátedra, Episodio X (tercera parte), «Derrota propia y derrota por bloqueo»; «Preferencia entre argumentos»}

\begin{center}
\small
\begin{tabular}{@{}c c c l c@{}}
\toprule
\textbf{Atacante} & \textbf{Sub-arg.\ de desacuerdo} & \textbf{Reglas} & \textbf{Comparación} & \textbf{Resultado} \\
\midrule
$B$ & $A_1$ & \{R1\} vs.\ \{R3, R5\} & sin relación entre R1 y R3, R5 & \textbf{bloqueo} \\
$C$ & $D$ (de $A_1$) & \{R6\} vs.\ \{R5\} & R6 $>$ R5: $C$ preferido & \textbf{propio} \\
$B$ & $A_2$ & \{R1\} vs.\ \{R4\} & R4 $>$ R1: $A_2$ preferido & no derrota \\
\bottomrule
\end{tabular}
\end{center}

\begin{itemize}
  \item $\langle B, \texttt{aire\_acond}\rangle$ es \textbf{derrotador por bloqueo} de $\langle A_1, \texttt{\nf aire\_acond}\rangle$: no son comparables.
  \item $\langle C, \texttt{\nf abrir\_ventana}\rangle$ es \textbf{derrotador propio} de $\langle A_1, \texttt{\nf aire\_acond}\rangle$: es preferido a $D$, el sub-argumento atacado.
  \item $\langle B, \texttt{aire\_acond}\rangle$ ataca a $\langle A_2, \texttt{\nf aire\_acond}\rangle$ pero \textbf{no es derrotador}: $A_2$ es preferido a $B$.
  \item La regla R2 aparece en el orden pero no influye: no está en ningún argumento.
\end{itemize}

\begin{center}
\begin{minipage}{\linewidth}
\centering
\begin{tikzpicture}
  \tri{-3}{0}{0.55}{1.2}{$C$}{\nf ab}{azul}
  \tri{0}{0}{1.1}{2.4}{}{\nf ac}{ambar}
  \node[font=\scriptsize\bfseries] at (0.55,0.25) {$A_1$};
  \tri{-0.2}{0.35}{0.45}{0.9}{$D$}{}{ambar}
  \node[font=\scriptsize\ttfamily, above] at (-0.2,1.17) {ab};
  \tri{3.2}{0}{0.55}{1.2}{$B$}{ac}{azul}
  \tri{6.2}{0}{0.55}{1.2}{$A_2$}{\nf ac}{ambar}
  \draw[dpropia] (-2.6,0.8) -- (-0.45,1.4);
  \node[font=\scriptsize, text=red!70!black] at (-1.7,1.45) {propia};
  \draw[dbloqueo] (2.85,0.9) -- (0.35,2.05);
  \node[font=\scriptsize, text=red!70!black] at (2.2,1.75) {bloqueo};
  \draw[ataque, {Stealth[length=2mm]}-] (5.8,0.6) -- (3.6,0.6);
  \node[font=\scriptsize, text=gris] at (4.7,0.85) {no derrota};
  \node[font=\scriptsize] at (-1.6,0.35) {R6 $>$ R5};
  \node[font=\scriptsize] at (4.7,0.3) {R4 $>$ R1};
\end{tikzpicture}
\captionof{figure}{Derrotas sobre los argumentos para \texttt{\nf aire\_acond}: $C$ derrota a $A_1$ (propia, en $D$), $B$ y $A_1$ se bloquean, y $B$ no derrota a $A_2$.}
\end{minipage}
\end{center}

\begin{extras}
\begin{itemize}
  \item \textbf{Ataques en sentido contrario}, que se usan en los ejercicios 6 y 8:\catedra{Episodio X (tercera parte)}{Derrota propia y derrota por bloqueo}
  \begin{itemize}
    \item $A_1$ ataca a $B$ y, como no son comparables, también es su \textbf{derrotador por bloqueo}.
    \item $A_2$ ataca a $B$ y es su \textbf{derrotador propio} (R4 $>$ R1).
    \item $D$ ataca a $C$ pero no lo derrota ($C$ es preferido). $A_1$ no ataca a $C$: \texttt{\nf aire\_acond} y \texttt{\nf abrir\_ventana} junto con $\Pi_5$ no forman un conjunto contradictorio; el ataque es solo de su sub-argumento $D$.
  \end{itemize}
\end{itemize}
\end{extras}

%======================================================================
\section{Ejercicio 6: Líneas de argumentación}

% Colores de soporte e interferencia en las líneas.
\colorlet{soporte}{azul}
\colorlet{interf}{ambar}

%----------------------------------------------------------------------
\subsection{a)}

\begin{enunciado}{Consigna}
¿Qué es una línea de argumentación? ¿Qué condiciones deben cumplirse para que una línea de argumentación sea aceptable?
\end{enunciado}

\paragraph{Línea de argumentación.} Una línea de argumentación para $\langle A_0, h_0\rangle$ es una \textbf{secuencia finita de argumentos} $[A_0, A_1, A_2, \dots]$ que comienza con $A_0$ y en la que, a partir del segundo, \textbf{cada argumento es un derrotador del anterior}.\catedra{Episodio X (tercera parte)}{Línea de argumentación para un argumento}
\begin{itemize}
  \item $A_1$ derrota a $A_0$, pero $A_2$ derrota a $A_1$, así que \textbf{defiende} a $A_0$; y así sucesivamente.
  \item \textbf{Argumentos de soporte} para $A_0$: los que lo defienden, $\{A_0, A_2, A_4, \dots\}$ (posiciones pares).
  \item \textbf{Argumentos de interferencia} para $A_0$: los que lo derrotan directa o indirectamente, $\{A_1, A_3, \dots\}$ (posiciones impares).
\end{itemize}

\begin{center}
\begin{minipage}{\linewidth}
\centering
\begin{tikzpicture}
  \foreach \i/\c in {0/soporte, 1/interf, 2/soporte, 3/interf, 4/soporte} {
    \tri{2*\i}{0}{0.5}{1.1}{$A_\i$}{$h_\i$}{\c}
  }
  \foreach \i in {1,...,4} {
    \draw[dpropia, draw=gris] (2*\i-0.45,0.55) -- (2*\i-1.55,0.55);
  }
  \node[font=\scriptsize, text=soporte, anchor=west] at (9,0.8) {soporte: $A_0, A_2, A_4$};
  \node[font=\scriptsize, text=interf, anchor=west] at (9,0.3) {interferencia: $A_1, A_3$};
\end{tikzpicture}
\captionof{figure}{Una línea de argumentación para $A_0$: cada argumento derrota al anterior (la flecha va del derrotador al derrotado).}
\end{minipage}
\end{center}

\paragraph{Línea aceptable.} Una línea de argumentación para $A_0$ es \textbf{aceptable} si:\fuentes{Cátedra, Episodio X (tercera parte), «Línea de argumentación aceptable». Lo que evita cada condición: «(1) evitando líneas con ciclos (infinitos)», «(2) evitando contradicciones», «(3) evitando doble bloqueo»}
\begin{enumerate}
  \item Ningún argumento $A_k$ \textbf{aparece más de una vez} en la línea, ni es un \textbf{sub-argumento de desacuerdo} de otro $A_i$ que aparece antes ($i < k$).
  \item El conjunto de argumentos de \textbf{soporte} para $A_0$ \textbf{no es contradictorio}, y tampoco el de \textbf{interferencia}. Un conjunto de argumentos $\{A_1, \dots, A_N\}$ es contradictorio si $\Pi \cup A_1 \cup \dots \cup A_N$ deriva un par de literales complementarios $p$ y ${\sim}p$.
  \item Si $A_k$ es un \textbf{derrotador por bloqueo} de $A_{k-1}$, entonces $A_{k+1}$ debe ser un \textbf{derrotador propio} de $A_k$: no puede haber dos derrotadores por bloqueo consecutivos.
\end{enumerate}

\begin{center}
\small
\begin{tabularx}{\linewidth}{@{}l X X@{}}
\toprule
\textbf{Condición} & \textbf{Exige} & \textbf{Evita} \\
\midrule
(1) Sin repetición & ningún argumento repetido ni sub-argumento de desacuerdo de uno anterior & líneas con ciclos (infinitas) \\
(2) Consistencia & soporte e interferencia no contradictorios por separado & que un bando use argumentos que se contradicen entre sí \\
(3) Sin doble bloqueo & después de un bloqueo, una derrota propia & dos derrotadores por bloqueo seguidos \\
\bottomrule
\end{tabularx}
\end{center}

\begin{center}
\begin{minipage}{\linewidth}
\centering
\begin{tikzpicture}
  % (1) ciclo
  \tri{0}{0}{0.45}{1}{$A_0$}{$h_0$}{soporte}
  \tri{1.5}{0}{0.45}{1}{$A_1$}{$h_1$}{interf}
  \tri{3}{0}{0.45}{1}{$A_2$}{$h_2$}{soporte}
  \tri{4.5}{0}{0.45}{1}{$A_1$}{$h_1$}{interf}
  \draw[dpropia, draw=gris] (1.15,0.5) -- (0.35,0.5);
  \draw[dpropia, draw=gris] (2.65,0.5) -- (1.85,0.5);
  \draw[dpropia, draw=gris] (4.15,0.5) -- (3.35,0.5);
  \node[font=\Large, text=red!70!black] at (4.5,0.45) {\texttimes};
  \node[font=\scriptsize, align=center] at (2.25,-0.5) {(1) $A_1$ repetido};
  % (3) doble bloqueo
  \tri{7}{0}{0.45}{1}{$A_0$}{a}{soporte}
  \tri{8.5}{0}{0.45}{1}{$A_1$}{\nf a}{interf}
  \tri{10}{0}{0.45}{1}{$A_2$}{a}{soporte}
  \draw[dbloqueo] (8.15,0.5) -- (7.35,0.5);
  \draw[dbloqueo] (9.65,0.5) -- (8.85,0.5);
  \node[font=\Large, text=red!70!black] at (10,0.45) {\texttimes};
  \node[font=\scriptsize, align=center] at (8.5,-0.5) {(3) dos bloqueos seguidos};
\end{tikzpicture}
\captionof{figure}{Dos líneas no aceptables: el argumento marcado con \texttimes{} no puede agregarse.}
\end{minipage}
\end{center}

\begin{extras}
\begin{itemize}
  \item \textbf{Por qué hace falta la condición (1)}: en el PLR Merval, $B$ (\texttt{buen\_precio(skynet)}) y $M$ (\texttt{\nf buen\_precio(skynet)}) se bloquean mutuamente, y sin restricciones la línea $[B, M, B, M, \dots]$ sería infinita.\fuentes{Cátedra, Episodio X (tercera parte), «¿ Argumentación infinita ?»; «Línea de argumentación aceptable»}
  \item \textbf{Por qué alcanza con no agregar el segundo bloqueo} (3): el bloqueo es simétrico; si $A_1$ bloquea a $A_0$, $A_0$ también bloquea a $A_1$, así que $A_1$ ya está bloqueado, se agregue o no $A_2$.\catedra{Episodio X (tercera parte)}{(3) evitando doble bloqueo}
  \item En un árbol de dialéctica, \textbf{todo camino de la raíz a una hoja es una línea de argumentación}.\catedra{Episodio X (tercera parte)}{Líneas de argumentación en un árbol de dialéctica}
\end{itemize}
\end{extras}

%----------------------------------------------------------------------
\subsection{b)}

\begin{enunciado}{Consigna}
Considerando el programa $\mathcal{P}_5$ del ejercicio anterior, determine todas las líneas de argumentación aceptables para los argumentos cuya conclusión sea el literal aire\_acond o el literal \nf aire\_acond.
\end{enunciado}

Programa $\mathcal{P}_5$, con los nombres de reglas del ejercicio 5 y el orden R2 $>$ R4, R4 $>$ R1, R6 $>$ R5:
\programaPcinco
\reglasPcinco

\paragraph{Argumentos} (del ejercicio 5):
\begin{itemize}
  \item $\langle A_1, \texttt{\nf aire\_acond}\rangle$, $A_1 = \{$R3, R5$\}$, con sub-argumento $\langle D, \texttt{abrir\_ventana}\rangle$, $D = \{$R5$\}$;
  \item $\langle A_2, \texttt{\nf aire\_acond}\rangle$, $A_2 = \{$R4$\}$;
  \item $\langle B, \texttt{aire\_acond}\rangle$, $B = \{$R1$\}$;
  \item $\langle C, \texttt{\nf abrir\_ventana}\rangle$, $C = \{$R6$\}$.
\end{itemize}

\paragraph{Derrotas} (ejercicio 5.b y sus EXTRAS):
\begin{center}
\small
\begin{tabular}{@{}l l l@{}}
\toprule
\textbf{Argumento} & \textbf{Sus derrotadores} & \textbf{Tipo} \\
\midrule
$A_1$ & $B$ & bloqueo \\
 & $C$ (en el sub-argumento $D$) & propio \\
$A_2$ & --- & \\
$B$ & $A_1$ & bloqueo \\
 & $A_2$ & propio \\
$C$ & --- ($D$ lo ataca, pero $C$ es preferido) & \\
\bottomrule
\end{tabular}
\end{center}

\paragraph{Líneas aceptables.} Se construyen agregando derrotadores mientras se cumplan las tres condiciones:
\begin{center}
\small
\begin{tabularx}{\linewidth}{@{}l l l l X@{}}
\toprule
\textbf{Raíz} & \textbf{Línea} & \textbf{Soporte} & \textbf{Interferencia} & \textbf{Por qué termina} \\
\midrule
$A_1$ & $\Lambda_1 = [A_1, C]$ & $\{A_1\}$ & $\{C\}$ & $C$ no tiene derrotadores \\
 & $\Lambda_2 = [A_1, B, A_2]$ & $\{A_1, A_2\}$ & $\{B\}$ & $A_2$ no tiene derrotadores \\
\midrule
$A_2$ & $\Lambda_3 = [A_2]$ & $\{A_2\}$ & $\emptyset$ & $A_2$ no tiene derrotadores \\
\midrule
$B$ & $\Lambda_4 = [B, A_2]$ & $\{B\}$ & $\{A_2\}$ & $A_2$ no tiene derrotadores \\
 & $\Lambda_5 = [B, A_1, C]$ & $\{B, C\}$ & $\{A_1\}$ & $C$ no tiene derrotadores \\
\bottomrule
\end{tabularx}
\end{center}

Por qué son aceptables:
\begin{itemize}
  \item \textbf{(1)}: en ninguna se repite un argumento, y ninguno es sub-argumento de desacuerdo de uno anterior.
  \item \textbf{(2)}: ningún conjunto de soporte o de interferencia es contradictorio. Los únicos con dos argumentos son $\{A_1, A_2\}$ (ambos concluyen \texttt{\nf aire\_acond}) y $\{B, C\}$ (\texttt{aire\_acond} y \texttt{\nf abrir\_ventana}); junto con $\Pi_5$ no derivan ningún par $p$, ${\sim}p$.
  \item \textbf{(3)}: en $\Lambda_2$ y $\Lambda_5$, al bloqueo del segundo argumento le sigue una \textbf{derrota propia} ($A_2$ sobre $B$; $C$ sobre $A_1$ en $D$).
\end{itemize}

Líneas descartadas:
\begin{itemize}
  \item $[A_1, B, A_1]$ y $[B, A_1, B]$: violan \textbf{(1)} (argumento repetido) y \textbf{(3)} (dos bloqueos seguidos).
\end{itemize}

\begin{aviso}{Líneas más cortas}
Los prefijos de una línea aceptable también son líneas aceptables (p.~ej.\ $[A_1]$, $[A_1, B]$, $[B]$, $[B, A_1]$). Arriba se muestran solo las \textbf{maximales}, que no se pueden extender: son las que aparecen como caminos de la raíz a una hoja en los árboles de dialéctica.
\end{aviso}

\begin{center}
\begin{minipage}{\linewidth}
\centering
\begin{tikzpicture}[lb/.style={font=\scriptsize\bfseries, anchor=east},
                    et/.style={font=\tiny, text=red!70!black, above=1pt}]
  \def\d{2.3}
  % L1
  \node[lb] at (-0.7,0.5) {$\Lambda_1$};
  \tri{0}{0}{0.45}{1}{$A_1$}{\nf ac}{soporte}
  \tri{\d}{0}{0.45}{1}{$C$}{\nf ab}{interf}
  \draw[dpropia] (\d-0.4,0.4) -- (0.4,0.4);
  \node[et] at (\d/2,0.4) {propia (en $D$)};
  % L2
  \node[lb] at (4.3,0.5) {$\Lambda_2$};
  \tri{5}{0}{0.45}{1}{$A_1$}{\nf ac}{soporte}
  \tri{5+\d}{0}{0.45}{1}{$B$}{ac}{interf}
  \tri{5+2*\d}{0}{0.45}{1}{$A_2$}{\nf ac}{soporte}
  \draw[dbloqueo] (5+\d-0.4,0.4) -- (5.4,0.4);
  \node[et] at (5+\d/2,0.4) {bloqueo};
  \draw[dpropia] (5+2*\d-0.4,0.4) -- (5+\d+0.4,0.4);
  \node[et] at (5+1.5*\d,0.4) {propia};
  % L3
  \node[lb] at (-0.7,-1.6) {$\Lambda_3$};
  \tri{0}{-2.1}{0.45}{1}{$A_2$}{\nf ac}{soporte}
  % L4
  \node[lb] at (1.6,-1.6) {$\Lambda_4$};
  \tri{2.3}{-2.1}{0.45}{1}{$B$}{ac}{soporte}
  \tri{2.3+\d}{-2.1}{0.45}{1}{$A_2$}{\nf ac}{interf}
  \draw[dpropia] (2.3+\d-0.4,-1.7) -- (2.7,-1.7);
  \node[et] at (2.3+\d/2,-1.7) {propia};
  % L5
  \node[lb] at (6.3,-1.6) {$\Lambda_5$};
  \tri{7}{-2.1}{0.45}{1}{$B$}{ac}{soporte}
  \tri{7+\d}{-2.1}{0.45}{1}{$A_1$}{\nf ac}{interf}
  \tri{7+2*\d}{-2.1}{0.45}{1}{$C$}{\nf ab}{soporte}
  \draw[dbloqueo] (7+\d-0.4,-1.7) -- (7.4,-1.7);
  \node[et] at (7+\d/2,-1.7) {bloqueo};
  \draw[dpropia] (7+2*\d-0.4,-1.7) -- (7+\d+0.4,-1.7);
  \node[et] at (7+1.5*\d,-1.7) {propia (en $D$)};
  % leyenda
  \node[font=\scriptsize, text=soporte, anchor=west] at (0,-3) {soporte};
  \node[font=\scriptsize, text=interf, anchor=west] at (1.5,-3) {interferencia};
  \node[font=\scriptsize, anchor=west] at (3.5,-3) {(\texttt{ac} = \texttt{aire\_acond}, \texttt{ab} = \texttt{abrir\_ventana})};
\end{tikzpicture}
\captionof{figure}{Las cinco líneas de argumentación aceptables maximales para los argumentos de \texttt{aire\_acond} y \texttt{\nf aire\_acond}.}
\end{minipage}
\end{center}

%======================================================================
\section{Ejercicio 7}

% Sub-argumento dentro de un argumento (relleno más oscuro).
\newcommand{\subtri}[7]{%
  \draw[thick, draw=#7, fill=#7!60] (#1-#3,#2) -- (#1,#2+#4) -- (#1+#3,#2) -- cycle;
  \node[font=\scriptsize\bfseries, text=white] at (#1,#2+0.3*#4) {#5};
  \node[font=\tiny\ttfamily, above=-1pt] at (#1,#2+#4) {#6};}
% Nombre de un argumento debajo de su triángulo.
\newcommand{\nomarg}[2]{\node[font=\scriptsize\bfseries, below] at (#1,0) {#2};}
\newcommand{\cruz}[2]{\node[font=\huge, text=red!70!black] at (#1,#2) {\texttimes};}

\begin{enunciado}{Consigna}
Indique si las siguientes líneas de argumentación son aceptables. Justifique adecuadamente en cada caso apelando a las condiciones enunciadas en el ejercicio 6.a.

(Siguen cuatro líneas, (a)--(d), dadas como figura; abajo se redibujan una por una.)
\end{enunciado}

\begin{aviso}{Cómo se lee la figura}
\begin{itemize}
  \item Cada triángulo es un argumento, con su conclusión arriba; un triángulo interior es un sub-argumento. Dos triángulos con el mismo nombre son el mismo argumento.
  \item Cada flecha es una derrota y apunta al \textbf{sub-argumento de desacuerdo} del derrotado (al vértice, si ataca la conclusión).
  \item La figura \textbf{no indica} si las derrotas son propias o por bloqueo, así que la condición (3) solo puede evaluarse de forma condicional.
\end{itemize}
\end{aviso}

Se aplican las tres condiciones del ejercicio 6.a. La que decide en los cuatro casos es la (1): el argumento que se agrega no puede estar \textbf{repetido} ni ser el \textbf{sub-argumento de desacuerdo} de uno anterior.\fuentes{Cátedra, Episodio X (tercera parte), «Línea de argumentación aceptable»; «(1) evitando líneas con ciclos (infinitos)». Las líneas (a) y (b) son las dos de esta última diapositiva}

%----------------------------------------------------------------------
\needspace{5cm}
\paragraph{(a) $[A_0, X, A_1, X]$: no es aceptable.}

\begin{center}
\begin{tikzpicture}
  \foreach \x/\n/\h/\c in {0/A_0/h0/soporte, 2.2/X/h1/interf, 4.4/A_1/h2/soporte, 6.6/X/h1/interf} {
    \tri{\x}{0}{0.6}{1.3}{}{\h}{\c}
    \subtri{\x}{0.1}{0.3}{0.6}{}{}{\c}
    \nomarg{\x}{$\n$}
  }
  \foreach \x in {2.2, 4.4, 6.6} {
    \draw[dpropia, draw=gris] (\x-0.05,1.3) -- (\x-2.2+0.2,0.42);
  }
  \cruz{6.6}{0.55}
\end{tikzpicture}
\captionof{figure}{Línea (a): $X$ aparece dos veces.}
\end{center}

\begin{itemize}
  \item $X$ aparece \textbf{dos veces} (posiciones 2 y 4): viola la condición \textbf{(1)}.
\end{itemize}

%----------------------------------------------------------------------
\paragraph{(b) $[A_0, A_1, X]$, con $A_1$ atacando en $\langle X, h_4\rangle$: no es aceptable.}

\begin{center}
\begin{tikzpicture}
  \tri{0}{0}{0.6}{1.3}{}{h0}{soporte}
  \subtri{0}{0.1}{0.32}{0.65}{X}{h4}{soporte}
  \nomarg{0}{$A_0$}
  \tri{2.2}{0}{0.6}{1.3}{}{h1}{interf}
  \subtri{2.2}{0.1}{0.3}{0.6}{}{}{interf}
  \nomarg{2.2}{$A_1$}
  \subtri{4.4}{0}{0.6}{1.3}{X}{}{soporte}
  \node[font=\scriptsize\ttfamily, above] at (4.4,1.3) {h4};
  \nomarg{4.4}{$X$}
  \draw[dpropia, draw=gris] (2.15,1.3) -- (0.24,0.4);
  \draw[dpropia, draw=gris] (4.35,1.3) -- (2.4,0.42);
  \cruz{4.4}{0.5}
\end{tikzpicture}
\captionof{figure}{Línea (b): $X$ es el sub-argumento de desacuerdo de $A_0$.}
\end{center}

\begin{itemize}
  \item $A_1$ derrota a $A_0$ atacando su sub-argumento $\langle X, h_4\rangle$: ese es el \textbf{sub-argumento de desacuerdo} de $A_0$.
  \item $X$, que ya fue derrotado por $A_1$, vuelve a entrar para derrotar a $A_1$: es sub-argumento de desacuerdo de un argumento anterior ($A_0$) y viola la condición \textbf{(1)}. Sería una argumentación circular: $X$ se defendería a sí mismo.
\end{itemize}

%----------------------------------------------------------------------
\paragraph{(c) $[A_0, A_1, X]$, con $A_1$ atacando la conclusión $h_0$: es aceptable (salvo doble bloqueo).}

\begin{center}
\begin{tikzpicture}
  \tri{0}{0}{0.6}{1.3}{}{h0}{soporte}
  \subtri{0}{0.1}{0.32}{0.65}{X}{h4}{soporte}
  \nomarg{0}{$A_0$}
  \tri{2.2}{0}{0.6}{1.3}{}{h1}{interf}
  \subtri{2.2}{0.1}{0.3}{0.6}{}{}{interf}
  \nomarg{2.2}{$A_1$}
  \subtri{4.4}{0}{0.6}{1.3}{X}{}{soporte}
  \node[font=\scriptsize\ttfamily, above] at (4.4,1.3) {h4};
  \nomarg{4.4}{$X$}
  \draw[dpropia, draw=gris] (2.1,1.3) -- (0.12,1.3);
  \draw[dpropia, draw=gris] (4.35,1.3) -- (2.4,0.42);
\end{tikzpicture}
\captionof{figure}{Línea (c): $X$ es sub-argumento de $A_0$, pero no el de desacuerdo.}
\end{center}

\begin{itemize}
  \item \textbf{(1) se cumple}: no hay repetidos, y $A_1$ ataca a $A_0$ en su conclusión, así que el sub-argumento de desacuerdo es $A_0$ completo. $X$ es sub-argumento de $A_0$, pero \textbf{no el de desacuerdo}.
  \item \textbf{(2) se cumple}:
  \begin{itemize}
    \item Soporte $\{A_0, X\}$: como $X \subseteq A_0$, $\Pi \cup A_0 \cup X = \Pi \cup A_0$, que no es contradictorio porque $A_0$ es un argumento.\catedra{Episodio X, parte 2}{Argumento}
    \item Interferencia $\{A_1\}$: un solo argumento, tampoco es contradictorio.
  \end{itemize}
  \item \textbf{(3)} depende del tipo de derrota: se viola solo si $A_1$ bloquea a $A_0$ \textbf{y} $X$ bloquea a $A_1$.
\end{itemize}

%----------------------------------------------------------------------
\paragraph{(d) $[A_0, A_1, X]$, con $A_1$ atacando en $\langle S, h_5\rangle$: es aceptable (salvo doble bloqueo).}

\begin{center}
\begin{tikzpicture}
  \tri{0}{0}{0.9}{1.6}{}{h0}{soporte}
  \subtri{-0.33}{0.1}{0.28}{0.6}{X}{h4}{soporte}
  \subtri{0.33}{0.1}{0.28}{0.6}{S}{h5}{soporte}
  \nomarg{0}{$A_0$}
  \tri{2.4}{0}{0.6}{1.3}{}{h1}{interf}
  \subtri{2.4}{0.1}{0.3}{0.6}{}{}{interf}
  \nomarg{2.4}{$A_1$}
  \subtri{4.6}{0}{0.6}{1.3}{X}{}{soporte}
  \node[font=\scriptsize\ttfamily, above] at (4.6,1.3) {h4};
  \nomarg{4.6}{$X$}
  \draw[dpropia, draw=gris] (2.35,1.3) -- (0.58,0.3);
  \draw[dpropia, draw=gris] (4.55,1.3) -- (2.6,0.42);
\end{tikzpicture}
\captionof{figure}{Línea (d): $A_1$ ataca a $A_0$ en $S$; $X$ es otro sub-argumento de $A_0$.}
\end{center}

\begin{itemize}
  \item \textbf{(1) se cumple}: el sub-argumento de desacuerdo de $A_0$ es $\langle S, h_5\rangle$ ($S$ es el nombre que se le da acá al sub-argumento para $h_5$), no $X$.
  \item \textbf{(2)} y \textbf{(3)}: igual que en (c), porque el soporte vuelve a ser $\{A_0, X\}$ con $X \subseteq A_0$ y los tipos de derrota tampoco se indican.
\end{itemize}

\paragraph{Resumen.}
\begin{center}
\small
\begin{tabularx}{\linewidth}{@{}l >{\raggedright\arraybackslash}X l l l@{}}
\toprule
\textbf{Línea} & \textbf{(1)} & \textbf{(2)} & \textbf{(3)} & \textbf{¿Aceptable?} \\
\midrule
(a) & no: $X$ repetido & --- & --- & \textbf{no} \\
(b) & no: $X$ es sub-arg.\ de desacuerdo de $A_0$ & --- & --- & \textbf{no} \\
(c) & sí: el desacuerdo es en $h_0$ & sí & si no hay doble bloqueo & \textbf{sí} \\
(d) & sí: el desacuerdo es en $h_5$ & sí & si no hay doble bloqueo & \textbf{sí} \\
\bottomrule
\end{tabularx}
\end{center}

\begin{extras}
\begin{itemize}
  \item \textbf{Ejemplo de la cátedra del caso (b)}, en el programa de \texttt{surf}: $\langle L, \texttt{\nf nice}\rangle$ derrota por bloqueo a $\langle A, \texttt{surf}\rangle$ en el sub-argumento de desacuerdo $\langle N, \texttt{nice}\rangle$. Aunque $N$ derrota por bloqueo a $L$, no puede agregarse después de $L$, porque es el sub-argumento de desacuerdo de $A$.\catedra{Episodio X (tercera parte)}{Verificando que la línea de argumentación sea aceptable}
\end{itemize}
\end{extras}

%======================================================================
\section{Ejercicio 8: Análisis dialéctico y respuesta a consultas}

% Nodos de árboles de dialéctica: U (azul) y D (ámbar). La flecha va del hijo (derrotador) al padre.
\tikzset{
  nodo/.style={draw, thick, rounded corners=2pt, align=center, font=\scriptsize, inner sep=2pt, minimum width=9mm, minimum height=6mm},
  nU/.style={nodo, draw=azul, fill=fondoazul, label={[font=\scriptsize\bfseries, text=azul]right:U}},
  nD/.style={nodo, draw=ambar, fill=fondoambar, label={[font=\scriptsize\bfseries, text=ambar]right:D}},
  podado/.style={nodo, draw=gris, dashed, text=gris},
  arbol/.style={level distance=12mm, sibling distance=20mm,
    edge from parent/.style={draw=gris, thick, {Stealth[length=2mm]}-}},
}

%----------------------------------------------------------------------
\subsection{a)}

\begin{enunciado}{Consigna}
¿Qué es un árbol dialéctico?
\end{enunciado}

\paragraph{Árbol de dialéctica para $\langle A_0, h_0\rangle$.}
\begin{itemize}
  \item La \textbf{raíz} se etiqueta con $\langle A_0, h_0\rangle$.\catedra{Episodio X (tercera parte)}{Definición de Árbol de dialéctica}
  \item Para un nodo $\langle A_N, h_N\rangle$, con camino desde la raíz $[A_0, \dots, A_N]$: cada \textbf{derrotador} $\langle B_i, q_i\rangle$ de $A_N$ tal que $[A_0, \dots, A_N, B_i]$ es una \textbf{línea de argumentación aceptable} es un \textbf{hijo} de $A_N$.
  \item Si $A_N$ no tiene ningún derrotador en esas condiciones, es una \textbf{hoja}.
\end{itemize}

\paragraph{Marcado.} Para saber si la raíz resulta derrotada, DeLP marca los nodos:\catedra{Episodio X (tercera parte)}{Árbol de dialéctica marcado}
\begin{itemize}
  \item \textbf{U} (\textit{undefeated}): toda \textbf{hoja}, y todo nodo interno cuyos hijos están \textbf{todos marcados D}.
  \item \textbf{D} (\textit{defeated}): todo nodo interno con \textbf{al menos un hijo marcado U}.
\end{itemize}

\begin{center}
\begin{minipage}{\linewidth}
\centering
\begin{tikzpicture}[arbol, sibling distance=30mm]
  \node[nU] {$A$\\\texttt{surf}}
    child {node[nD] {$L$\\\texttt{\nf nice}} child {node[nU] {$R$\\\texttt{\nf rain}}}}
    child {node[nD] {$B$\\\texttt{busy}} child {node[nU] {$G$\\\texttt{\nf yard\_work}}}};
\end{tikzpicture}
\captionof{figure}{Árbol de dialéctica marcado de la cátedra para \texttt{surf}: las hojas $R$ y $G$ son U, así que $L$ y $B$ son D y la raíz $A$ es U. La flecha va del derrotador al derrotado.}
\end{minipage}
\end{center}

\begin{extras}
\begin{itemize}
  \item En ese ejemplo, $L$ y $B$ derrotan a $A$ en sus sub-argumentos de desacuerdo $\langle N, \texttt{nice}\rangle$ y $\langle T, \texttt{\nf busy}\rangle$, y $R$ y $G$ son derrotadores propios de $L$ y $B$ en sus sub-argumentos para \texttt{rain} y \texttt{yard\_work}. La respuesta para \texttt{surf} es \textbf{yes}.\fuentes{Cátedra, Episodio X (tercera parte), «Marcado completo del árbol»; «Respuesta: sí, disfrutemos de un día de surf»}
\end{itemize}
\end{extras}

%----------------------------------------------------------------------
\subsection{b)}

\begin{enunciado}{Consigna}
Si la raíz de un árbol dialéctico está marcada como U, ¿es posible asegurar que todos los argumentos que defienden a la raíz van a estar marcados U, y todos aquellos que atacan (directa o indirectamente) a la raíz van a estar marcados D? Muestre un ejemplo que justifique su respuesta.
\end{enunciado}

\textbf{No.}
\begin{itemize}
  \item Que la raíz sea U solo exige que \textbf{sus hijos} sean todos D, y a cada hijo D le alcanza con \textbf{un} hijo U. Los demás nodos de esos subárboles pueden tener cualquier marca.
  \item Ejemplo (árbol de la cátedra): la raíz $B$ es U, pero el argumento de soporte $E_2$ (nivel 2, defiende a $B$) está marcado \textbf{D}, y el de interferencia $F$ (nivel 3) está marcado \textbf{U}.\catedra{Episodio X (tercera parte)}{Conclusión garantizada}
\end{itemize}

\begin{center}
\begin{minipage}{\linewidth}
\centering
\begin{tikzpicture}[arbol]
  \node[nU] {$B$\\\texttt{c}}
    child {node[nD] {$C_1$} child {node[nU] {$E_1$}}}
    child {node[nD] {$C_2$}
      child {node[nD, line width=1.6pt] {$E_2$} child {node[nU, line width=1.6pt] {$F$}}}
      child {node[nU] {$E_3$}}};
  \node[font=\scriptsize, anchor=west, align=left] at (3.2,-1.2) {nivel 1: interferencia};
  \node[font=\scriptsize, anchor=west, align=left] at (3.2,-2.4) {nivel 2: soporte};
  \node[font=\scriptsize, anchor=west, align=left] at (3.2,-3.6) {nivel 3: interferencia};
\end{tikzpicture}
\captionof{figure}{Raíz U con un argumento de soporte D ($E_2$) y uno de interferencia U ($F$), resaltados.}
\end{minipage}
\end{center}

%----------------------------------------------------------------------
\subsection{c)}

\begin{enunciado}{Consigna}
¿Es necesario analizar todas las ramas de un árbol dialéctico para conocer el marcado de los nodos? Muestre un ejemplo que justifique su respuesta.
\end{enunciado}

\textbf{No}, no siempre. Se sigue de la regla de marcado:\catedra{Episodio X (tercera parte)}{Árbol de dialéctica marcado}
\begin{itemize}
  \item Apenas un nodo tiene \textbf{un hijo U}, ya es \textbf{D}: los subárboles de sus otros hijos no cambian su marca y pueden no explorarse (\textbf{poda}).
  \item Para que un nodo sea \textbf{U}, en cambio, hay que comprobar que \textbf{todos} sus hijos sean D.
\end{itemize}

Ejemplo, con $\mathcal{P}_5$: el árbol para $\langle B, \texttt{aire\_acond}\rangle$ tiene dos hijos, $A_2$ y $A_1$ (ejercicio 6.b). Si primero se explora $A_2$, que es una hoja (U), $B$ queda marcado D sin analizar el subárbol de $A_1$.

\begin{center}
\begin{minipage}{\linewidth}
\centering
\begin{tikzpicture}[arbol, sibling distance=26mm]
  \node[nD] {$B$\\\texttt{a\_a}}
    child {node[nU] {$A_2$\\\texttt{\nf a\_a}}}
    child {node[podado] {$A_1$\\\texttt{\nf a\_a}} child {node[podado] {$C$\\\texttt{\nf ab\_v}}} edge from parent[dashed]};
  \node[font=\scriptsize, text=gris, anchor=west] at (2.1,-1.8) {no hace falta explorarlo};
\end{tikzpicture}
\captionof{figure}{Poda: como $A_2$ es U, $B$ es D sin importar la rama de $A_1$ (punteada). \texttt{a\_a} $=$ \texttt{aire\_acond}, \texttt{ab\_v} $=$ \texttt{abrir\_ventana}.}
\end{minipage}
\end{center}

\begin{extras}
\begin{itemize}
  \item La cátedra pide a veces los árboles ``\textbf{sin podas}'', es decir, completos, aunque no hagan falta para el marcado.\fuentes{Cátedra, «Resolución con DeLP» (ejemplo resuelto)}
\end{itemize}
\end{extras}

%----------------------------------------------------------------------
\subsection{d)}

\begin{enunciado}{Consigna}
¿Cuándo se dice que un literal está garantizado? ¿Qué relación existe entre la garantía de un literal y la de su complemento?
\end{enunciado}

\begin{itemize}
  \item Un literal $c$ está \textbf{garantizado} por un programa DeLP si existe \textbf{al menos un árbol de dialéctica} para un argumento que sustenta $c$ cuya \textbf{raíz está marcada U}.\catedra{Episodio X (tercera parte)}{Conclusión garantizada}
  \item Por eso, ante una consulta $c$, DeLP analiza los argumentos para $c$ y los argumentos para ${\sim}c$.
  \item \textbf{Relación con el complemento}:
  \begin{itemize}
    \item $c$ y ${\sim}c$ \textbf{no pueden estar garantizados a la vez}.\fuentesweb{No está en la cátedra. García \& Simari, \textit{Defeasible Logic Programming: An Argumentative Approach}, TPLP 4(1--2), 2004 (referencia [2] del Episodio X, tercera parte)}
    \item Puede que \textbf{ninguno} esté garantizado (situación de indecisión).\catedra{Episodio X (tercera parte)}{Situación de indecisión}
  \end{itemize}
\end{itemize}

\begin{center}
\small
\begin{tabular}{@{}l l l@{}}
\toprule
\textbf{$c$} & \textbf{${\sim}c$} & \textbf{¿Posible?} \\
\midrule
garantizado & no garantizado & sí \\
no garantizado & garantizado & sí \\
no garantizado & no garantizado & sí (indecisión) \\
garantizado & garantizado & \textbf{no} \\
\bottomrule
\end{tabular}
\end{center}

\begin{extras}
\begin{itemize}
  \item En los ejemplos de la cátedra, un literal puede tener varios argumentos y varios árboles: alcanza con que \textbf{uno} tenga la raíz U, aunque los otros tengan la raíz D.\catedra{Episodio X (tercera parte)}{Conclusión garantizada}
\end{itemize}
\end{extras}

%----------------------------------------------------------------------
\subsection{e)}

\begin{enunciado}{Consigna}
¿Cuáles son las posibles respuestas de un intérprete DeLP?
\end{enunciado}

Para una consulta $c$ a un programa $\mathcal{P}$:\fuentes{Cátedra, Episodio X (tercera parte), «Respuesta a una consulta». Misma definición en Cátedra, «Resolución con DeLP»}
\begin{itemize}
  \item \textbf{yes}: $c$ está garantizado por $\mathcal{P}$.
  \item \textbf{no}: el \textbf{complemento} de $c$ está garantizado por $\mathcal{P}$ (complemento$(c) = {\sim}c$ y complemento$({\sim}c) = c$).
  \item \textbf{undecided}: $c$ pertenece a $\mathcal{P}$, pero ni $c$ ni su complemento están garantizados.
  \item \textbf{unknown}: $c$ \textbf{no pertenece} a $\mathcal{P}$, es decir, ni $c$ ni su complemento aparecen en alguna regla de $\mathcal{P}$.
\end{itemize}

\begin{center}
\begin{minipage}{\linewidth}
\centering
\begin{tikzpicture}[node distance=5mm and 7mm, every node/.style={font=\small}]
  \node[cajaazul] (q) {¿$c$ aparece en $\mathcal{P}$\\(como $c$ o ${\sim}c$)?};
  \node[caja, right=of q] (unk) {\textbf{unknown}};
  \node[cajaazul, below=of q] (g) {¿$c$ garantizado?};
  \node[caja, right=of g] (yes) {\textbf{yes}};
  \node[cajaazul, below=of g] (gn) {¿${\sim}c$ garantizado?};
  \node[caja, right=of gn] (no) {\textbf{no}};
  \node[caja, below=of gn] (und) {\textbf{undecided}};
  \draw[flecha] (q) -- node[above, font=\scriptsize] {no} (unk);
  \draw[flecha] (q) -- node[left, font=\scriptsize] {sí} (g);
  \draw[flecha] (g) -- node[above, font=\scriptsize] {sí} (yes);
  \draw[flecha] (g) -- node[left, font=\scriptsize] {no} (gn);
  \draw[flecha] (gn) -- node[above, font=\scriptsize] {sí} (no);
  \draw[flecha] (gn) -- node[left, font=\scriptsize] {no} (und);
\end{tikzpicture}
\captionof{figure}{Respuesta del intérprete para una consulta $c$.}
\end{minipage}
\end{center}

%----------------------------------------------------------------------
\subsection{f)}

\begin{enunciado}{Consigna}
Determine qué literales están garantizados según el programa $\mathcal{P}_5$, construyendo los árboles dialécticos correspondientes.
\end{enunciado}

Programa $\mathcal{P}_5$, con el orden R2 $>$ R4, R4 $>$ R1, R6 $>$ R5:
\programaPcinco
\reglasPcinco

\paragraph{Argumentos y derrotas.} Son los del ejercicio 6.b ($A_1$, $A_2$, $B$, $C$). Hay que agregar los dos que faltan:
\begin{itemize}
  \item $\langle D, \texttt{abrir\_ventana}\rangle$, $D = \{$R5$\}$ (sub-argumento de $A_1$). Su único derrotador es $C$, \textbf{propio} (R6 $>$ R5): \texttt{abrir\_ventana} y \texttt{\nf abrir\_ventana} están en desacuerdo.
  \item Los literales de $\Pi_5$ (\texttt{calor}, \texttt{caluroso}, \texttt{viento\_caluroso}, \texttt{ahorrar\_energia} y \texttt{viento}) tienen como argumento el \textbf{conjunto vacío}.\catedra{Episodio X, parte 2}{Argumento} No tienen derrotadores: un argumento para el complemento sería contradictorio con $\Pi_5$, así que no sería argumento.
  \item \texttt{sofocante} no tiene argumento (no se deriva; ver el aviso del ejercicio 5), y por eso \texttt{aire\_acond} tampoco tiene un argumento con R2.
\end{itemize}

\paragraph{Árboles.}

\begin{center}
\begin{minipage}{\linewidth}
\centering
% Argumento simple: \argt{x}{y}{nombre}{conclusión}{marca}{color}
\newcommand{\argt}[6]{%
  \tri{#1}{#2}{0.42}{0.95}{#3}{#4}{#6}
  \node[font=\scriptsize\bfseries, text=#6, anchor=west] at (#1+0.42,#2+0.3) {#5};}
% A1 con su sub-argumento D adentro: \argAuno{x}{y}{marca}{color}
\newcommand{\argAuno}[4]{%
  \tri{#1}{#2}{0.8}{1.6}{}{\nf a\_a}{#4}
  \subtri{#1-0.28}{#2+0.07}{0.32}{0.78}{D}{ab\_v}{#4}
  \node[font=\scriptsize\bfseries] at (#1+0.33,#2+0.25) {$A_1$};
  \node[font=\scriptsize\bfseries, text=#4, anchor=west] at (#1+0.75,#2+0.3) {#3};}
\begin{tikzpicture}[der/.style={-{Stealth[length=2mm]}, thick, draw=gris}]
  % (1) A1
  \argAuno{1.2}{0}{D}{ambar}
  \argt{-0.1}{-2.6}{$B$}{a\_a}{D}{ambar}
  \argt{2.5}{-2.6}{$C$}{\nf ab\_v}{U}{azul}
  \argt{-0.1}{-5.2}{$A_2$}{\nf a\_a}{U}{azul}
  \draw[der] (-0.1,-1.3) -- (1.45,0);
  \draw[der] (2.5,-1.3) -- (1.1,0.42);
  \draw[der] (-0.1,-3.9) -- (-0.1,-2.6);
  % (2) A2
  \argt{4.3}{0}{$A_2$}{\nf a\_a}{U}{azul}
  % (3) B
  \argt{7.3}{0}{$B$}{a\_a}{D}{ambar}
  \argt{6.0}{-2.6}{$A_2$}{\nf a\_a}{U}{azul}
  \argAuno{8.6}{-2.6}{D}{ambar}
  \argt{9.1}{-5.2}{$C$}{\nf ab\_v}{U}{azul}
  \draw[der] (6.0,-1.3) -- (7.15,0);
  \draw[der] (8.6,-0.7) -- (7.45,0);
  \draw[der] (9.1,-3.9) -- (8.5,-2.18);
  % (4) C
  \argt{10.8}{0}{$C$}{\nf ab\_v}{U}{azul}
  % (5) D
  \argt{12.8}{0}{$D$}{ab\_v}{D}{ambar}
  \argt{12.8}{-2.6}{$C$}{\nf ab\_v}{U}{azul}
  \draw[der] (12.8,-1.3) -- (12.8,0);
  \foreach \x/\t in {1.2/(1), 4.3/(2), 7.3/(3), 10.8/(4), 12.8/(5)} {
    \node[font=\scriptsize\bfseries] at (\x,2.3) {\t};
  }
\end{tikzpicture}
\captionof{figure}{Árboles de dialéctica marcados de $\mathcal{P}_5$ (azul: U; ámbar: D). \texttt{a\_a} $=$ \texttt{aire\_acond}, \texttt{ab\_v} $=$ \texttt{abrir\_ventana}. El triángulo interior de $A_1$ es su sub-argumento $\langle D, \texttt{ab\_v}\rangle$. La flecha va del derrotador al punto de ataque: al sub-argumento $D$, o a la base del triángulo si ataca la conclusión.}
\end{minipage}
\end{center}

\begin{enumerate}[label=(\arabic*)]
  \item $A_1$: hijos $B$ (bloqueo) y $C$ (propio, en $D$). Bajo $B$ va $A_2$ (propio); $A_1$ no puede volver a entrar (condición 1). $A_2$ y $C$ son hojas: U. $B$ queda D y $A_1$ queda \textbf{D}.
  \item $A_2$: no tiene derrotadores. Hoja: \textbf{U}.
  \item $B$: hijos $A_2$ (propio) y $A_1$ (bloqueo). Bajo $A_1$ va $C$ (propio, en $D$), sin doble bloqueo. $A_2$ y $C$ son U, $A_1$ es D y $B$ queda \textbf{D}.
  \item $C$: no tiene derrotadores ($D$ lo ataca, pero $C$ es preferido). Hoja: \textbf{U}.
  \item $D$: hijo $C$ (propio). $C$ es U, así que $D$ queda \textbf{D}.
\end{enumerate}

\paragraph{Literales garantizados.}
\begin{center}
\small
\begin{tabularx}{\linewidth}{@{}>{\raggedright\arraybackslash}p{4.6cm} l >{\raggedright\arraybackslash}X@{}}
\toprule
\textbf{Literal} & \textbf{¿Garantizado?} & \textbf{Por qué} \\
\midrule
\texttt{calor}, \texttt{caluroso}, \texttt{viento\_caluroso}, \texttt{ahorrar\_energia}, \texttt{viento} & sí & argumento $\emptyset$ sin derrotadores: árbol de un solo nodo, U \\
\texttt{\nf aire\_acond} & \textbf{sí} & árbol (2) con raíz U (aunque el árbol (1) tenga raíz D) \\
\texttt{\nf abrir\_ventana} & \textbf{sí} & árbol (4) con raíz U \\
\texttt{aire\_acond} & no & su único árbol (3) tiene raíz D \\
\texttt{abrir\_ventana} & no & su único árbol (5) tiene raíz D \\
\texttt{sofocante} & no & no tiene argumento \\
\bottomrule
\end{tabularx}
\end{center}

\begin{extras}
\begin{itemize}
  \item \textbf{Respuestas del intérprete} para estos literales, según las definiciones del inciso e):\catedra{Episodio X (tercera parte)}{Respuesta a una consulta}
  \begin{center}
  \small
  \begin{tabularx}{\linewidth}{@{}>{\raggedright\arraybackslash}p{5.2cm} >{\raggedright\arraybackslash}X@{}}
  \toprule
  \textbf{Consulta} & \textbf{Respuesta} \\
  \midrule
  literales de $\Pi_5$, \texttt{\nf aire\_acond}, \texttt{\nf abrir\_ventana} & yes \\
  \texttt{aire\_acond}, \texttt{abrir\_ventana} & no (su complemento está garantizado) \\
  \texttt{sofocante} & undecided (aparece en R2, pero ni él ni su complemento están garantizados) \\
  \bottomrule
  \end{tabularx}
  \end{center}
\end{itemize}
\end{extras}

%----------------------------------------------------------------------
\subsection{g)}

\begin{enunciado}{Consigna}
¿Puede aparecer el mismo argumento más de una vez en el mismo árbol dialéctico? En caso afirmativo, ¿dicho argumento siempre va a quedar marcado igual?
\end{enunciado}

\begin{itemize}
  \item \textbf{Sí}: un argumento no puede repetirse dentro de una \textbf{línea} (condición 1 del ejercicio 6.a), pero sí en \textbf{líneas diferentes} del mismo árbol.\catedra{Episodio X (tercera parte)}{Argumentos repetidos en un árbol de dialéctica}
  \item \textbf{No} siempre queda marcado igual: la marca depende del subárbol que tiene debajo, y en cada línea los derrotadores admitidos pueden ser distintos.
\end{itemize}

Ejemplo de la cátedra:\catedra{Episodio X (tercera parte)}{Argumentos en diferentes partes de un árbol de dialéctica}
\begin{itemize}
  \item $\Pi = \{x, y, z, u\}$; preferencia $(b \reb e) > ({\sim}b \reb y)$.
  \item $A = \{a \reb d;\ d \reb x\}$, $C_2 = \{{\sim}a \reb {\sim}b, e;\ {\sim}b \reb y;\ e \reb u\}$, $C_3 = \{{\sim}d \reb {\sim}b;\ {\sim}b \reb y\}$, $X = \{b \reb e;\ e \reb u\}$, $C_4 = \{{\sim}e \reb z\}$.
  \item $X$ aparece dos veces, las dos como argumento de soporte para $A$:
  \begin{itemize}
    \item En $[A, C_2, X]$ es \textbf{U}: $C_4$ no puede derrotarlo ahí, porque $C_2$ y $C_4$ son de interferencia y contradictorios ($e$ y ${\sim}e$), lo que viola la condición 2.
    \item En $[A, C_3, X, C_4]$ es \textbf{D}: esa línea sí es aceptable, y $C_4$ (hoja, U) lo derrota.
  \end{itemize}
\end{itemize}

\begin{center}
\begin{minipage}{\linewidth}
\centering
\begin{tikzpicture}[arbol, sibling distance=28mm]
  \node[nD] {$A$\\$a$}
    child {node[nD] {$C_2$\\${\sim}a$} child {node[nU, line width=1.6pt] {$X$\\$b$}}}
    child {node[nU] {$C_3$\\${\sim}d$} child {node[nD, line width=1.6pt] {$X$\\$b$} child {node[nU] {$C_4$\\${\sim}e$}}}};
\end{tikzpicture}
\captionof{figure}{El mismo argumento $X$ (resaltado) aparece dos veces en el árbol, con marcas distintas.}
\end{minipage}
\end{center}

\begin{extras}
\begin{itemize}
  \item En otro programa de la cátedra, $X$ aparece \textbf{tres veces} en el mismo árbol, siempre U pero en niveles distintos: dos veces como soporte y una como interferencia, porque derrota a tres argumentos diferentes.\catedra{Episodio X (tercera parte)}{Argumentos repetidos en un árbol de dialéctica}
  \item También puede cambiar el \textbf{rol}: en un tercer ejemplo, $X$ aparece una vez como soporte (U) y otra como interferencia (D).\catedra{Episodio X (tercera parte)}{Argumentos repetidos en diferentes líneas y con diferente rol y marca}
\end{itemize}
\end{extras}

%----------------------------------------------------------------------
\subsection{h)}

\begin{enunciado}{Consigna}
Además de para determinar la respuesta del intérprete ¿para qué más nos pueden servir los árboles dialécticos?
\end{enunciado}

\begin{itemize}
  \item Como \textbf{explicación} de la respuesta: el árbol muestra qué argumentos sustentan la conclusión, cuáles la atacan, quién derrota a quién y por qué la raíz queda U o D.\fuentes{Cátedra, Episodio X (tercera parte), «Respuesta a una consulta» («Los árboles de dialéctica son, además, una explicación para las respuestas»); «Respuesta: sí, disfrutemos de un día de surf»}
\end{itemize}

\begin{extras}
\begin{itemize}
  \item La cátedra lo cuenta entre las características de DeLP: \textit{es posible dar una explicación detallada de las respuestas}. Una de sus referencias trata el tema: García \& Simari, \textit{Defeasible Logic Programming: DeLP-servers, Contextual Queries, and Explanations for Answers} (2014).\catedra{Episodio X (tercera parte)}{Comentarios finales}
\end{itemize}
\end{extras}

\end{document}
